% Chapter: first-principles (density-functional) inputs of the device model.
% Writer-local environment (plain style, as the report's other note environments).
\newenvironment{dftplain}{\par\medskip\noindent\textit{In plain words.}\ }{\par\medskip}

\chapter{First-principles inputs: density-functional calculations of \texorpdfstring{\InO}{In2O3} and IWO films}
\label{ch:dft}

The device model of the previous chapters needs to know how the electronic properties of an \IWO{} film change when
the film becomes only one or two nanometres thick: by how much the band gap opens, how the opening is shared between
the conduction and the valence band, and how heavy the electrons become. It also needs to know what a tungsten atom
does inside indium oxide. None of these quantities has been measured for the target films. Until this chapter they
were taken over from published calculations on pure \InO{} and carried the provenance label
\prov{DFT\_PURE\_IN2O3\_PROXY} with an assigned uncertainty of $\pm50$\,\% (\cref{sec:par-dec,sec:par-mstar}).

This chapter describes a first-principles study carried out to replace those proxies by values computed in one
consistent setup, and to add what the literature does not give for \IWO. Because the method is far from the device
physics of the other chapters, the chapter is written so that it can be read without prior knowledge of
density-functional theory (DFT). It proceeds in the order in which the work was done: why the calculations are
needed (\cref{sec:dft-motivation}); the vocabulary (\cref{sec:dft-glossary}); how a DFT calculation works
(\cref{sec:dft-theory}); how the material is defined for the computer (\cref{sec:dft-material}); every calculation
with its settings and result (\cref{sec:dft-calcs}); how the calculations were run on the CERN batch system and on
GPUs (\cref{sec:dft-compute}); how the numbers and figures are obtained from the output files
(\cref{sec:dft-analysis}); and what changes in the device model (\cref{sec:dft-pending}). Each idea is first
explained in plain words and then stated precisely. All numbers are read from the output files of the calculations
by scripts; the dated log of every value is \file{qe_workflow/cern_htcondor/results/RESULTS_LOG.md}.

Method references that are not bundled with the package (marked by the absence of a page number) are cited for the
method as a whole; no page-specific statement is taken from them. Statements about Quantum ESPRESSO, Lin et al.\ and
Si et al.\ carry the printed page of the bundled text.

% =====================================================================================================================
\section{Why first-principles calculations are needed}
\label{sec:dft-motivation}

\begin{dftplain}
A transistor channel that is only 2~nm thick is a few atomic layers of oxide. In such a thin film the electrons are
squeezed between the two surfaces, and quantum mechanics then raises the lowest energy they can have (the band gap
opens) and changes how easily they move (their effective mass grows). The device simulator cannot work this out by
itself; it must be told. The numbers can come from measurements on the films, which do not exist, or from a
calculation that starts from the atoms alone. Such a calculation is called ``first-principles'' or ``ab initio'',
because it uses no adjustable parameters fitted to the material, only the laws of quantum mechanics and the positions
of the atoms.
\end{dftplain}

The thickness dependence enters the device model through two laws, the confinement shift of the conduction-band edge
$\dEc(t) = 0.9205\,t^{-1.3815}$~eV ($t$ in nm, \cref{eq:dec-law}), fitted to the slab-minus-bulk gaps of
\citet[p.~21542]{Lin2022} and the conduction-band shift of \citet[p.~504]{Si2021}, and an increment of the effective mass
based on the slab masses of \citet[p.~21540]{Lin2022}. Three weaknesses motivated new calculations:
\begin{enumerate}
\item The proxies describe pure \InO; the channel is W-doped, and the role of W was unknown.
\item The published values come from different codes and settings, and their uncertainty could only be assigned
($\pm50$\,\%), not estimated.
\item The semi-local functional used in all of them is known to fail for band gaps (\cref{sec:dft-xc}); whether it is
adequate for the \emph{change} of the gap with thickness had not been tested.
\end{enumerate}

The study therefore asks five questions:
\begin{description}[leftmargin=1.2em]
\item[Q1] Does an independent calculation reproduce the published confinement of a 1~nm film?
\item[Q2] What are the confinement shift and the mass at 2~nm, the thinnest device?
\item[Q3] How is the gap opening shared between conduction and valence band ($\dEc$ versus $\Delta E_{\mathrm{v}}$)?
\item[Q4] Is the semi-local functional adequate for the shift, as opposed to the gap itself?
\item[Q5] Where does W sit in the lattice, which valence does it take, and what happens to its electrons?
\end{description}
\Cref{fig:dft-pipeline} shows how the answers travel from the atoms to the device model.

\begin{figure}[htbp]\centering
\begin{tikzpicture}[node distance=4mm and 5mm, every node/.style={font=\footnotesize},
  box/.style={draw, rounded corners=2pt, align=center, minimum height=11mm, text width=24mm, fill=lightbg},
  arr/.style={-{Stealth[length=2mm]}, thick}]
\node[box] (s) {Atomic structure\\(crystal, film, W)\\\cref{sec:dft-material}};
\node[box, right=of s] (d) {DFT calculation\\(Quantum ESPRESSO)\\\cref{sec:dft-theory,sec:dft-calcs}};
\node[box, right=of d] (o) {Energies, bands,\\electrostatic potential\\(output files)};
\node[box, right=of o] (n) {$\Delta E_{\mathrm{g}}$, $\dEc$, $\mstar$,\\W site, valence\\\cref{sec:dft-analysis}};
\node[box, below=7mm of n] (t) {Thickness laws of\\the TCAD model\\\cref{sec:dft-pending}};
\node[box, left=of t, text width=50mm] (c) {CERN batch system (CPU, GPU), workflow\\\cref{sec:dft-compute}};
\draw[arr] (s) -- (d); \draw[arr] (d) -- (o); \draw[arr] (o) -- (n); \draw[arr] (n) -- (t);
\draw[arr, dashed] (c.north) -- ++(0,3.5mm) -| (d.south);
\end{tikzpicture}
\caption{From atoms to device inputs: the path of this chapter. The structure is defined, the DFT calculation
produces energies, bands and potentials, analysis scripts turn them into the quantities the device model uses, and
these replace the literature proxies of the thickness laws.}
\label{fig:dft-pipeline}
\end{figure}

What the calculations can and cannot deliver must be stated at the start. They describe ideal crystalline films with
hydrogen-terminated surfaces in vacuum, whereas the device channels are sputtered, largely amorphous or nanocrystalline
films on \AlO{} with an air-exposed top. The results are therefore \emph{computed} inputs that describe trends and
give physical grounding; they are not \emph{validated} against measurements of the target films (\cref{sec:dft-limits}).

% =====================================================================================================================
\section{Terms used in this chapter}
\label{sec:dft-glossary}

\Cref{tab:dft-glossary} explains every technical term of the chapter in plain words, roughly in the order in which the
terms appear. \Cref{tab:dft-units} lists the units: DFT codes work in ``atomic'' units (Rydberg, bohr), the device
model in eV and nm.

\begin{small}
\begin{longtable}{@{}>{\raggedright\arraybackslash}p{3.3cm}>{\raggedright\arraybackslash}p{9.0cm}>{\raggedright\arraybackslash}p{1.6cm}@{}}
\caption{Glossary of the terms used in this chapter.}\label{tab:dft-glossary}\\
\toprule
Term & Meaning in plain words & Section \\
\midrule
\endfirsthead
\multicolumn{3}{@{}l}{\textit{\Cref{tab:dft-glossary}, continued}}\\
\toprule
Term & Meaning in plain words & Section \\
\midrule
\endhead
\bottomrule
\endfoot
\multicolumn{3}{@{}l}{\textit{Physics of electrons in solids}}\\
Band, band structure & In a crystal the allowed electron energies form continuous ranges called bands. The band
structure is the plot of these energies against the electron's wave vector $k$. & \ref{sec:dft-bloch} \\
Valence band (VB), conduction band (CB) & The highest band filled with electrons, and the lowest empty band above it.
Current in an n-type oxide is carried by electrons in the CB. & \ref{sec:dft-bloch} \\
VBM, CBM & Valence-band maximum and conduction-band minimum: the top of the VB and the bottom of the CB. & \ref{sec:dft-readout} \\
Band gap $\Eg$ & The energy distance CBM $-$ VBM. ``Fundamental'' gap: smallest distance anywhere in $k$;
``direct'' gap: distance at one $k$ (here $\Gamma$). & \ref{sec:dft-readout} \\
Confinement shift $\Delta E_{\mathrm{g}}(t)$ & How much the gap of a film of thickness $t$ exceeds the gap of the thick
(bulk) crystal, because the film squeezes the electrons. & \ref{sec:dft-readout} \\
$\dEc$, $\Delta E_{\mathrm{v}}$ & The parts of $\Delta E_{\mathrm{g}}$ that come from the CB moving up and the VB moving
down. & \ref{sec:dft-readout} \\
Effective mass $\mstar$ & How heavy an electron appears when it moves in the crystal; read from the curvature of the
CB near its minimum. Given in units of the free-electron mass $m_0$. & \ref{sec:dft-readout} \\
Non-parabolicity $\alpha$ & How much the CB deviates from a parabola away from its minimum (the mass grows with energy). & \ref{sec:dft-readout} \\
Wave vector $k$, $\Gamma$ & A label of an electron state in a crystal (its momentum divided by $\hbar$). $\Gamma$ is
$k = 0$. & \ref{sec:dft-bloch} \\
Brillouin zone (BZ) & The set of all distinct wave vectors of a crystal. Calculations sum over points of this zone. & \ref{sec:dft-bloch} \\
$k$-point grid & The finite set of $k$ points used for those sums, e.g.\ 3$\times$3$\times$3. Denser grids are more
accurate and more expensive. & \ref{sec:dft-bloch} \\
Fermi level $E_F$ & The energy up to which states are filled. In a metal it lies inside a band. & \ref{sec:dft-occ} \\
Density of states (DOS), projected DOS (PDOS) & How many electron states exist per energy interval; the PDOS
splits this into contributions of each atom and orbital (e.g.\ the W 5d orbitals). & \ref{sec:dft-readout} \\
Vacuum level $E_{\mathrm{vac}}$ & The energy of an electron at rest just outside the material; the common reference that
lets band edges of different films be compared. & \ref{sec:dft-readout} \\
Electron affinity (EA), ionization potential (IP) & $E_{\mathrm{vac}} - $ CBM and $E_{\mathrm{vac}} - $ VBM: the depths
of the band edges below the vacuum level. & \ref{sec:dft-readout} \\
Spin polarization, magnetic moment & Whether spin-up and spin-down electrons are distributed differently; the
difference, in Bohr magnetons ($\mu_{\mathrm{B}}$), is the magnetic moment. & \ref{sec:dft-calc-w} \\
\midrule
\multicolumn{3}{@{}l}{\textit{The DFT method}}\\
Density-functional theory (DFT) & A way to compute the ground state of many interacting electrons from the electron
density alone, instead of the many-electron wave function. & \ref{sec:dft-ks} \\
Electron density $n(\mathbf{r})$ & The number of electrons per volume at each point in space. & \ref{sec:dft-ks} \\
Kohn--Sham equations, orbitals, eigenvalues & The one-electron equations that DFT actually solves. Their solutions
(orbitals) and energies (eigenvalues) are used as the band structure. & \ref{sec:dft-ks} \\
Exchange--correlation (xc) functional & The one approximate ingredient of DFT: the quantum-mechanical part of the
electron--electron interaction. & \ref{sec:dft-xc} \\
PBE & The semi-local (generalized-gradient) xc functional used for all main results. Good for structures and
energy differences, but gives band gaps that are far too small. & \ref{sec:dft-xc} \\
Hybrid functional, HSE06 & An xc functional that mixes in 25\,\% of exact (Fock) exchange at short range. Gives much
better gaps, at much higher cost. Used here as a check. & \ref{sec:dft-xc} \\
Plane waves, cutoff (\texttt{ecutwfc}, \texttt{ecutrho}) & The orbitals are written as sums of plane waves; the cutoff
is the highest kinetic energy kept. A higher cutoff means more plane waves and more accuracy. & \ref{sec:dft-pw} \\
Pseudopotential, PAW, norm-conserving (NC) & Replacing the tightly bound core electrons and the steep nuclear potential
by a smooth effective potential, so that few plane waves suffice. PAW (projector augmented wave) and NC are two kinds. & \ref{sec:dft-pp} \\
Self-consistent field (SCF) & The iterative loop: guess the density, compute the potential, solve for the orbitals,
build a new density, repeat until input and output agree. & \ref{sec:dft-scf} \\
Mixing, \texttt{mixing\_beta}, local-TF & How the next input density is formed from previous ones; a small
\texttt{mixing\_beta} is a cautious step; local-TF is a variant designed for slabs. & \ref{sec:dft-scf} \\
Convergence threshold & The tolerance at which an iteration is declared finished (e.g.\ \texttt{conv\_thr} for the SCF). & \ref{sec:dft-scf} \\
Occupations: fixed, smearing & Insulators: every band either full or empty. Metals: occupations smeared over a small
energy width so that sums over $k$ converge. & \ref{sec:dft-occ} \\
Total energy, force, stress, pressure & The ground-state energy of the whole cell; its derivatives with respect to atom
positions (forces) and cell shape (stress; pressure is its mean). & \ref{sec:dft-relax} \\
Relaxation (\texttt{relax}, \texttt{vc-relax}), BFGS & Moving the atoms (and in vc-relax also the cell) step by step
downhill in energy until the forces vanish. BFGS is the step algorithm. & \ref{sec:dft-relax} \\
Non-SCF (``bands'') calculation & A run that keeps the converged density fixed and only computes orbitals at
chosen $k$ points, e.g.\ along lines for a band-structure plot. & \ref{sec:dft-readout} \\
Planar / macroscopic average & The electrostatic potential averaged over each plane parallel to the film, and then
over one atomic-layer period, to remove atomic wiggles. & \ref{sec:dft-readout} \\
\midrule
\multicolumn{3}{@{}l}{\textit{Describing the material}}\\
Unit cell, conventional / primitive cell, supercell & The repeating box of a crystal. The conventional cubic cell of
\InO{} has 80 atoms; the smallest (primitive) cell has 40. A supercell is several cells stacked. & \ref{sec:dft-bulkstruct} \\
Bixbyite, space group Ia$\bar{3}$ & The crystal structure of \InO{} and its symmetry group (number 206). & \ref{sec:dft-bulkstruct} \\
Wyckoff site (8b, 24d, 48e) & A class of symmetry-equivalent positions in the cell; the number is how many there are.
In \InO{} indium occupies 8b and 24d, oxygen 48e. & \ref{sec:dft-bulkstruct} \\
Symmetry operation & A rotation, reflection or shift that maps the crystal onto itself; their number is a quick check
that a structure is correct. & \ref{sec:dft-bulkstruct} \\
Slab, vacuum layer & A film model: a few atomic layers, repeated periodically with empty space (vacuum) between the
copies so that they do not interact. & \ref{sec:dft-slabstruct} \\
Passivation & Saturating the broken bonds of the surface oxygen atoms with hydrogen so that no spurious surface states
appear in the gap. & \ref{sec:dft-slabstruct} \\
Surface dipole & An electric dipole across the film when its two surfaces differ; it would tilt the vacuum potential.
The slabs here have equivalent surfaces and no dipole. & \ref{sec:dft-slabstruct} \\
Substitution, cation fraction & Replacing one In atom by W; the cation fraction W/(W+In) is the W content. & \ref{sec:dft-wstruct} \\
\midrule
\multicolumn{3}{@{}l}{\textit{Computing}}\\
Quantum ESPRESSO (QE), \texttt{pw.x}, \texttt{pp.x}, \texttt{projwfc.x} & The open-source DFT package used; \texttt{pw.x}
solves the equations, \texttt{pp.x} extracts potentials, \texttt{projwfc.x} computes the PDOS. & \ref{sec:dft-calcs} \\
MPI rank, pool (\texttt{-nk}) & The parallel processes of one calculation, and groups of them that each handle a subset
of the $k$ points. & \ref{sec:dft-driver} \\
CPU, GPU, core-hour & Processor types; GPUs (graphics processors) run QE several times faster for large cells. A
core-hour is one CPU core used for one hour. & \ref{sec:dft-gpu} \\
HTCondor, job, flavour, wall time & CERN's batch system; a job is one queued calculation; the flavour sets its maximum
run time (wall time). & \ref{sec:dft-cern} \\
lxplus, EOS, Kerberos ticket & CERN's login computers, its disk storage, and the time-limited login credential. & \ref{sec:dft-cern} \\
Container (Apptainer) & A packaged software environment; used to run NVIDIA's GPU build of QE. & \ref{sec:dft-gpu} \\
Checkpoint, continuation & A saved copy of a job's results after each step, and a resubmission that restarts an
unfinished relaxation from its last geometry. & \ref{sec:dft-flow} \\
\end{longtable}
\end{small}

\begin{table}[htbp]\centering\small
\caption{Units used by the DFT code and their conversion (CODATA 2018 values, as in the analysis scripts).}
\label{tab:dft-units}
\begin{tabular}{@{}lll@{}}
\toprule
Unit & Value & Used for \\
\midrule
Rydberg (Ry) & 1~Ry $= 13.605693$~eV & energies, cutoffs, convergence thresholds \\
bohr & 1~bohr $= 0.529177$~\AA & lengths, forces (Ry/bohr) \\
kbar & 1~kbar $= 0.1$~GPa & stress and pressure \\
$\hbar^2/2m_0$ & $3.80998$~eV\,\AA$^2$ & effective-mass fits \\
$\mu_{\mathrm{B}}$ & Bohr magneton & magnetic moments \\
\bottomrule
\end{tabular}
\end{table}

% =====================================================================================================================
\section{How a density-functional calculation works}
\label{sec:dft-theory}

This section explains the method from the beginning. The textbook treatment is \citet{Martin2004}; the
implementation used is Quantum ESPRESSO \citep{Giannozzi2009,Giannozzi2017}.

\subsection{The problem: many interacting electrons}
\label{sec:dft-problem}

\begin{dftplain}
All properties of a solid follow from its electrons, which move around the much heavier nuclei and repel each other.
Quantum mechanics describes them exactly by one equation, but that equation couples every electron to every other one.
For the 176-atom film of this study, which has 1328 electrons, it cannot be solved directly. DFT is the standard way
around this.
\end{dftplain}

Because nuclei are thousands of times heavier than electrons, the electrons can be treated as moving in the field of
fixed nuclei at positions $\{\mathbf{R}_I\}$ \citep{Born1927}. The electrons then obey
\begin{equation}
\Big[\sum_i -\frac{\hbar^2}{2m_0}\nabla_i^2 + \sum_i V_{\mathrm{ext}}(\mathbf{r}_i)
+ \frac{1}{2}\sum_{i\neq j}\frac{e^2}{4\pi\varepsilon_0|\mathbf{r}_i-\mathbf{r}_j|}\Big]\Psi = E\,\Psi ,
\label{eq:dft-manybody}
\end{equation}
where $V_{\mathrm{ext}}$ is the attraction of the nuclei and $\Psi(\mathbf{r}_1,\ldots,\mathbf{r}_N)$ depends on all $N$
electron coordinates at once. Storing $\Psi$ on even a coarse grid is impossible for $N$ in the hundreds.

\subsection{The density instead of the wave function: Kohn--Sham equations}
\label{sec:dft-ks}

\begin{dftplain}
DFT rests on a theorem: the ground-state energy is fully determined by the electron density $n(\mathbf{r})$, a function
of only three coordinates. Kohn and Sham then showed how to get that density by solving one-electron equations: each
electron moves in an average potential created by the nuclei and by all the other electrons. Everything that is hard
about the interaction is collected in one term, the exchange--correlation energy, which has to be approximated.
\end{dftplain}

\citet{Hohenberg1964} proved that the ground-state energy is a functional $E[n]$ of the density and is minimal for the
true density. \citet{Kohn1965} wrote the density as that of fictitious non-interacting electrons in orbitals $\psi_i$,
\begin{equation}
n(\mathbf{r}) = \sum_i f_i\,|\psi_i(\mathbf{r})|^2 ,
\label{eq:dft-density}
\end{equation}
with occupations $f_i$, and obtained the Kohn--Sham equations
\begin{equation}
\Big[-\frac{\hbar^2}{2m_0}\nabla^2 + V_{\mathrm{ext}}(\mathbf{r}) + V_{\mathrm{H}}[n](\mathbf{r})
+ V_{\mathrm{xc}}[n](\mathbf{r})\Big]\psi_i(\mathbf{r}) = \varepsilon_i\,\psi_i(\mathbf{r}) .
\label{eq:dft-ks}
\end{equation}
$V_{\mathrm{H}}$ is the classical (Hartree) repulsion of the density with itself and
$V_{\mathrm{xc}} = \delta E_{\mathrm{xc}}/\delta n$ the exchange--correlation potential. The total energy is
\begin{equation}
E[n] = T_{\mathrm{s}}[n] + \int V_{\mathrm{ext}}\,n\,d\mathbf{r} + E_{\mathrm{H}}[n] + E_{\mathrm{xc}}[n] + E_{\mathrm{II}} ,
\label{eq:dft-energy}
\end{equation}
with $T_{\mathrm{s}}$ the kinetic energy of the orbitals and $E_{\mathrm{II}}$ the repulsion between nuclei. Because
$V_{\mathrm{H}}$ and $V_{\mathrm{xc}}$ depend on $n$, which depends on the orbitals, \cref{eq:dft-ks} must be solved
self-consistently (\cref{sec:dft-scf}). The eigenvalues $\varepsilon_i$ are used as the band energies; strictly they
are auxiliary quantities, which is the origin of the band-gap problem discussed next.

\subsection{The approximation: exchange--correlation functionals}
\label{sec:dft-xc}

\begin{dftplain}
The exchange--correlation term is the only approximation in DFT, and the choice of approximation decides what the
calculation gets right. The common choice (PBE) gives good crystal structures and energy differences, but it places the
conduction band too low: band gaps come out about half or less of the measured value. More expensive ``hybrid''
functionals (HSE06) correct most of this. This study uses PBE for everything and checks with HSE06 that the
\emph{change} of the gap with film thickness is reliable, even though the gap itself is not.
\end{dftplain}

All main calculations use the generalized-gradient approximation of \citet{Perdew1996} (PBE), in which $E_{\mathrm{xc}}$
depends on $n$ and $|\nabla n|$ at each point. Semi-local functionals underestimate band gaps, partly because the
Kohn--Sham gap lacks the discontinuity of the exact exchange--correlation potential at integer electron number
\citep{PerdewLevy1983,ShamSchluter1983}. For \InO{} the PBE gap of \citet[p.~21542]{Lin2022} is 0.94~eV, against
experimental values of 2.7--3.25~eV quoted on the same page; the device model uses 3.05~eV. Lattice constants are
typically overestimated by 1--2\,\% with PBE (\cref{sec:dft-calc-relax}).

The hybrid functional HSE06 \citep{Heyd2003,Heyd2006,Krukau2006} splits the Coulomb interaction into a short-range and a
long-range part with a screening parameter $\omega$ and replaces a quarter of the short-range PBE exchange by exact
(Fock) exchange:
\begin{equation}
E_{\mathrm{xc}}^{\mathrm{HSE}} = \tfrac{1}{4}E_{\mathrm{x}}^{\mathrm{HF,SR}}(\omega) + \tfrac{3}{4}E_{\mathrm{x}}^{\mathrm{PBE,SR}}(\omega)
+ E_{\mathrm{x}}^{\mathrm{PBE,LR}}(\omega) + E_{\mathrm{c}}^{\mathrm{PBE}} .
\label{eq:dft-hse}
\end{equation}
The output of the HSE06 runs confirms the exact-exchange fraction of 0.25. Exact exchange couples every occupied orbital
at $\mathbf{k}$ with every orbital at $\mathbf{k}' = \mathbf{k}+\mathbf{q}$; Quantum ESPRESSO evaluates these terms with
fast Fourier transforms, at a cost at least an order of magnitude above a semi-local calculation, and allows the sum
over $\mathbf{q}$ to run over a coarser auxiliary grid whose adequacy must be checked
\citep[pp.~16--17]{Giannozzi2009}. The implementation in the versions used applies the exchange operator in the
adaptively compressed form of \citet{LinL2016}.

\subsection{Crystals: Bloch states, wave vectors and k-point sampling}
\label{sec:dft-bloch}

\begin{dftplain}
A crystal repeats itself, so the computer only needs one repeating box (the unit cell). The price is that each electron
state carries a label $k$, its wave vector, and the properties of the crystal are averages over all $k$ in a finite
region called the Brillouin zone. In practice the average is taken over a grid of $k$ points. Plotting the energies
against $k$ along lines through this zone gives the familiar band-structure diagram.
\end{dftplain}

In a periodic potential every orbital can be written as a plane wave times a cell-periodic function
\citep{Bloch1929}, $\psi_{n\mathbf{k}}(\mathbf{r}) = e^{i\mathbf{k}\cdot\mathbf{r}}\,u_{n\mathbf{k}}(\mathbf{r})$, with
band index $n$ and wave vector $\mathbf{k}$ in the first Brillouin zone. Densities and energies become averages over
$\mathbf{k}$, evaluated on Monkhorst--Pack grids \citep{Monkhorst1976}: an $N_1\times N_2\times N_3$ grid samples the
zone uniformly; for slabs the grid is $3\times3\times1$, because there is no dispersion across the vacuum. Quantum
ESPRESSO detects the crystal symmetries and restricts the sum to the irreducible wedge of the zone with appropriate
weights, then symmetrizes the density, forces and stress \citep[pp.~7, 16]{Giannozzi2009}. The bulk primitive cell has
4 irreducible points at $3\times3\times3$ and 8 at $4\times4\times4$; the 1~nm slab has 4 at $3\times3\times1$
(output files of the runs).

\subsection{Plane waves and the cutoff energy}
\label{sec:dft-pw}

\begin{dftplain}
To put the orbitals into a computer, each one is written as a sum of simple waves of different wavelengths. Short
wavelengths are needed to describe sharp features; the cutoff energy decides how short the shortest wave is. A higher
cutoff gives a more accurate answer and costs more. The cutoff is therefore chosen by testing: it is raised until the
results stop changing (\cref{sec:dft-calc-conv}).
\end{dftplain}

The periodic part of each orbital is expanded in plane waves with reciprocal-lattice vectors $\mathbf{G}$,
\begin{equation}
\psi_{n\mathbf{k}}(\mathbf{r}) = \sum_{\mathbf{G}} c_{n,\mathbf{k}+\mathbf{G}}\,e^{i(\mathbf{k}+\mathbf{G})\cdot\mathbf{r}},
\qquad \frac{\hbar^2}{2m_0}|\mathbf{k}+\mathbf{G}|^2 \le E_{\mathrm{cut}} ,
\label{eq:dft-pw}
\end{equation}
where $E_{\mathrm{cut}}$ is \texttt{ecutwfc}. In Rydberg units $\hbar^2/2m_0 = 1$, so the condition reads
$|\mathbf{k}+\mathbf{G}|^2 \le E_{\mathrm{cut}}$ in bohr$^{-2}$, and the number of plane waves is close to the number of
lattice points in a sphere of radius $\sqrt{E_{\mathrm{cut}}}$,
\begin{equation}
N_{\mathrm{PW}} \approx \frac{\Omega\,E_{\mathrm{cut}}^{3/2}}{6\pi^2} ,
\label{eq:dft-npw}
\end{equation}
with $\Omega$ the cell volume in bohr$^3$. For the relaxed bulk primitive cell ($\Omega = 3693.6$~bohr$^3$) and
71~Ry this gives 37\,315, against 37\,359 plane waves at $\Gamma$ in the output; for the relaxed 1~nm slab
($\Omega = 25\,140$~bohr$^3$) it gives 253\,982, against 253\,973. The density, a product of orbitals, needs a higher
cutoff \texttt{ecutrho}: for PAW and ultrasoft data sets typically 6--12 times \texttt{ecutwfc}, for norm-conserving
pseudopotentials 4 times \citep[p.~7]{Giannozzi2009}. This work uses 8 (568~Ry) with PAW and 4 (320~Ry) with the
norm-conserving set of the hybrid check. The density of the 1~nm slab is then represented by 5\,746\,731 plane waves on
a $150\times150\times500$ fast-Fourier-transform grid.

\subsection{Pseudopotentials and the projector augmented-wave method}
\label{sec:dft-pp}

\begin{dftplain}
Near a nucleus the orbitals oscillate rapidly, which would need an enormous number of plane waves. The core electrons
there do not take part in bonding. Pseudopotentials replace the nucleus and its core electrons by a smooth effective
potential that acts on the outer (valence) electrons only. The PAW method does this while keeping a way to reconstruct
the true orbitals near the nucleus, so it is accurate at a moderate cutoff.
\end{dftplain}

The atomic cores are described by the projector augmented-wave (PAW) method \citep{Bloechl1994,Kresse1999} with the
pslibrary 1.0.0 data sets \citep{DalCorso2014}; Quantum ESPRESSO supports norm-conserving, ultrasoft and PAW
descriptions \citep[pp.~5, 7]{Giannozzi2009}. \Cref{tab:dft-pp} lists the data sets. The valence configuration decides
how many electrons the calculation contains and therefore how many bands are occupied. For the hybrid-functional check,
optimized norm-conserving Vanderbilt pseudopotentials \citep{Hamann2013} of the SG15 library \citep{Schlipf2015} were
used instead (\cref{sec:dft-hse}).

\begin{table}[htbp]\centering\small
\caption{Pseudopotential data sets (PAW, pslibrary 1.0.0). Valence charge and suggested cutoffs from the file headers;
the calculations use 71/568~Ry, above every suggestion.}
\label{tab:dft-pp}
\begin{tabular}{@{}llrr@{}}
\toprule
Element & File & Valence electrons & Suggested cutoffs (Ry) \\
\midrule
In & \file{In.pbe-dn-kjpaw_psl.1.0.0.UPF} & 13 (4d$^{10}$5s$^2$5p$^1$) & 41 / 212 \\
O  & \file{O.pbe-n-kjpaw_psl.1.0.0.UPF}  & 6 (2s$^2$2p$^4$) & 47 / 323 \\
W  & \file{W.pbe-spn-kjpaw_psl.1.0.0.UPF} & 14 (5s$^2$5p$^6$5d$^4$6s$^2$) & 50 / 475 \\
H  & \file{H.pbe-kjpaw_psl.1.0.0.UPF}    & 1 & 46 / 221 \\
\bottomrule
\end{tabular}
\end{table}

Counting electrons is a useful check of every structure. The bulk primitive cell In$_{16}$O$_{24}$ has
$16\times13 + 24\times6 = 352$ electrons and therefore 176 occupied bands; the 1~nm slab In$_{24}$O$_{48}$H$_{24}$ has
624 (312 occupied bands); the 2~nm slab In$_{56}$O$_{96}$H$_{24}$ has 1328 (664). A W atom carries one electron more
than In in these data sets (14 against 13), but the chemically relevant count is different: In gives three electrons
to oxygen (In$^{3+}$), whereas W in its highest oxidation state gives six (W$^{6+}$), so each W substituted for In adds
three electrons that oxygen cannot take. The doped 80-atom cell has 705 electrons, an odd number, and is a metal in
the calculation.

\subsection{The self-consistent field loop}
\label{sec:dft-scf}

\begin{dftplain}
The potential in the Kohn--Sham equation depends on the density, and the density depends on the solutions. The
calculation therefore guesses a density, computes the potential, solves the equations, builds the new density, and
repeats. Each round is one ``SCF iteration''. To avoid oscillations, the new guess is a careful blend of the previous
ones (``mixing''). The loop stops when the density no longer changes by more than a tiny tolerance.
\end{dftplain}

\Cref{fig:dft-scf} shows the loop. Self-consistency is recast as a fixed-point problem $x = F[x]$ for the Fourier
components of the density, and Quantum ESPRESSO solves it with the modified Broyden method of \citet{Johnson1988}: the
next input density is the optimal combination of the inputs of a few previous iterations, plus a step along the
residual \citep[pp.~7, 13--14]{Giannozzi2009}. The step can be preconditioned by an approximate model of electronic
screening: simple, Thomas--Fermi, or local Thomas--Fermi (local-TF) mixing \citep[p.~14]{Giannozzi2009}
\citep{Raczkowski2001}. Local-TF accounts for the strongly inhomogeneous density of a slab with vacuum, and is used for
all slabs here; \texttt{mixing\_beta} sets the size of the step (0.4 bulk, 0.2 slabs, 0.3 doped bulk).

The convergence measure is an estimate of the error of the total energy built from the density residual,
$\|\rho^{\mathrm{out}}-\rho^{\mathrm{in}}\|^2 \propto \sum_{\mathbf{G}} |\Delta\rho(\mathbf{G})|^2/G^2$
\citep[p.~14]{Giannozzi2009}, printed as ``estimated scf accuracy''. The loop stops when it falls below
\texttt{conv\_thr} $= 10^{-9}$~Ry for every calculation of this study. Within each iteration the Kohn--Sham equations
are solved iteratively for the lowest bands only, by the block Davidson method or by band-by-band conjugate-gradient
minimization \citep[pp.~14--15]{Giannozzi2009} \citep{Davidson1975,Payne1992}. Davidson is faster; conjugate gradient
(CG) is slower but more robust (\cref{sec:dft-calc-bands}).

\begin{figure}[htbp]\centering
\begin{tikzpicture}[node distance=5mm, every node/.style={font=\footnotesize},
  box/.style={draw, rounded corners=2pt, align=center, fill=lightbg, minimum height=9mm, text width=44mm},
  dec/.style={draw, diamond, aspect=2.6, align=center, inner sep=1pt, text width=26mm},
  arr/.style={-{Stealth[length=2mm]}, thick}]
\node[box] (g) {Start: density from overlapping atoms\\(or from the previous geometry)};
\node[box, below=of g] (v) {Potential $V_{\mathrm{ext}} + V_{\mathrm{H}}[n] + V_{\mathrm{xc}}[n]$};
\node[box, below=of v] (k) {Solve Kohn--Sham equations at every $k$\\(Davidson or CG), lowest \texttt{nbnd} bands};
\node[box, below=of k] (o) {Occupy bands (fixed or smearing),\\new density $n^{\mathrm{out}}$, total energy};
\node[dec, below=6mm of o] (c) {estimated error\\$<$ \texttt{conv\_thr}?};
\node[box, right=14mm of v, text width=34mm] (m) {Mix: $n^{\mathrm{in}}$ from\\previous $n^{\mathrm{in}}$, $n^{\mathrm{out}}$\\(Broyden, local-TF)};
\node[box, below=6mm of c] (e) {Converged: energies, eigenvalues,\\forces, stress, density saved};
\draw[arr] (g) -- (v); \draw[arr] (v) -- (k); \draw[arr] (k) -- (o); \draw[arr] (o) -- (c);
\draw[arr] (c.east) -| node[pos=0.25, above] {no: next iteration} (m.south);
\draw[arr] (m.west) -- (v.east);
\draw[arr] (c) -- node[right] {yes} (e);
\end{tikzpicture}
\caption{The self-consistent field (SCF) loop of a Kohn--Sham calculation, as implemented in Quantum ESPRESSO
\citep[pp.~13--15]{Giannozzi2009}. In this study the loop stops at an estimated energy error of $10^{-9}$~Ry; it
takes 25 iterations for the bulk cell and 47 for the relaxed 1~nm slab.}
\label{fig:dft-scf}
\end{figure}

\subsection{Insulators and metals: occupations}
\label{sec:dft-occ}

\begin{dftplain}
In an insulator every band is either completely full or completely empty, and the program simply fills the lowest ones.
In a metal a band is partly filled, and the boundary between filled and empty states (the Fermi level) cuts through it.
Sums over a finite $k$ grid then converge badly, so the occupations are smoothed over a small energy window.
\end{dftplain}

Pure \InO{} cells and slabs are insulators and use fixed occupations. The W-doped cells and slabs have three extra
electrons per W in the conduction band and are metallic in the calculation; they use the cold smearing of
\citet{Marzari1999} with a width of 0.01~Ry, one of the smearing schemes of Quantum ESPRESSO \citep[p.~7]{Giannozzi2009}.
The Fermi level $E_F$ is then determined by the electron count.

\subsection{Forces, stress and relaxation}
\label{sec:dft-relax}

\begin{dftplain}
Once the electrons are in their ground state, the program can compute the force on every atom and the stress on the
cell. If the forces are not zero, the atoms are not where they want to be. A relaxation moves them step by step
downhill in energy, recomputing the electrons at each step, until the forces are negligible. In a variable-cell
relaxation the size and shape of the box are adjusted as well.
\end{dftplain}

Forces are the derivatives of the total energy with respect to the atomic positions, obtained from the ground state by
the Hellmann--Feynman theorem \citep{Feynman1939}, and the stress tensor follows from the derivative with respect to a
strain of the cell \citep{Nielsen1985}; the pressure is minus one third of its trace. Quantum ESPRESSO computes both
and uses them for structural optimizations of atomic coordinates and cell \citep[p.~5]{Giannozzi2009}. The
optimization uses the Broyden--Fletcher--Goldfarb--Shanno (BFGS) quasi-Newton algorithm
\citep[p.~7]{Giannozzi2009}: each step moves along $-\mathbf{H}^{-1}\nabla E$, where the inverse
Hessian $\mathbf{H}^{-1}$ is built up from the change of the gradients between steps. Variable-cell runs (\texttt{vc-relax})
use the invariant formulation of \citet{Wentzcovitch1991}. For slabs only the in-plane cell vectors are allowed to change
(\texttt{cell\_dofree = '2Dxy'}), so that the vacuum stays fixed. A relaxation ends when the energy change between steps
and the largest force component fall below \texttt{etot\_conv\_thr} and \texttt{forc\_conv\_thr}, and, for vc-relax,
the pressure below \texttt{press\_conv\_thr} (values in \cref{tab:dft-settings}). At the end, pw.x repeats the SCF with
the plane-wave basis of the final cell, which removes the error of a basis set that was defined for the starting cell
(Pulay stress).

\subsection{What is read from a calculation}
\label{sec:dft-readout}

\begin{dftplain}
A finished calculation gives a list of energies at each $k$ point. The highest filled and lowest empty ones are the
band edges, and their difference is the gap. The curvature of the conduction band near its minimum gives the effective
mass. The average electrostatic potential far out in the vacuum is the vacuum level, which serves as a common zero, so
that band edges of different films can be compared. \Cref{fig:dft-levels} shows how these quantities are defined.
\end{dftplain}

\paragraph{Band edges and gaps.} pw.x prints the highest occupied and lowest unoccupied eigenvalue of every SCF; their
difference is the Kohn--Sham gap on the $k$ grid (which contains $\Gamma$). For the bulk the fundamental gap is taken
from a separate non-self-consistent (``bands'') run along a path through the Brillouin zone, which uses the converged
density and computes orbitals at the requested $k$ points only.

\paragraph{Effective mass.} Near the conduction-band minimum at $\Gamma$ the band is fitted with a parabola,
\begin{equation}
E(k) = E_{\mathrm{CBM}} + \frac{\hbar^2 k^2}{2\mstar} , \qquad
\frac{\mstar}{m_0} = \frac{\hbar^2/2m_0}{c_2} = \frac{3.80998~\mathrm{eV}\,\text{\AA}^{2}}{c_2} ,
\label{eq:dft-mass}
\end{equation}
where $c_2$ is the slope of $E - E_{\mathrm{CBM}}$ against $k^2$ in a least-squares fit over $|k| \le 0.05$~\AA$^{-1}$.
Over a wider range the band is described by the non-parabolic (Kane-type) dispersion \citep{Kane1957}
\begin{equation}
E\,(1 + \alpha E) = \frac{\hbar^2 k^2}{2\mstar} ,
\label{eq:dft-kane}
\end{equation}
which is linear in the two unknowns $1/\mstar$ and $\alpha$ for known $E$ and $k$ and is solved by linear least squares.

\paragraph{Density of states and its projections.} projwfc.x projects the orbitals on atomic orbitals (L\"owdin
projections, \citealp{Lowdin1950}; \citealp[p.~10]{Giannozzi2009}) and gives the density of states resolved by atom and
orbital, here used for the W 5d share of the states near $E_F$.

\paragraph{Vacuum level, EA and IP.} pp.x computes the electrostatic potential felt by an electron (bare ionic plus
Hartree; \texttt{plot\_num = 11}), and average.x forms its planar average over each plane parallel to the film,
\begin{equation}
\bar V(z) = \frac{1}{A}\iint_A V(x,y,z)\,dx\,dy ,
\label{eq:dft-planar}
\end{equation}
and the macroscopic average $\bar{\bar V}(z)$, a further running average over one structural period $L$ that removes the
atomic oscillations \citep{Baldereschi1988}\citep[p.~10]{Giannozzi2009}; here $L = a/4$, the distance between cation
layers along [001]. In the middle of the vacuum the potential is flat and defines $E_{\mathrm{vac}}$. Then
\begin{equation}
\mathrm{EA} = E_{\mathrm{vac}} - E_{\mathrm{CBM}}, \qquad \mathrm{IP} = E_{\mathrm{vac}} - E_{\mathrm{VBM}}, \qquad
\Eg = \mathrm{IP} - \mathrm{EA} .
\label{eq:dft-eaip}
\end{equation}
Because both surfaces of the slabs are equivalent by symmetry, the potential is the same in the vacuum on both sides and
no dipole correction \citep{Neugebauer1992} is needed. The confinement shift and its partition between two films $t_1$
and $t_2$ follow as
\begin{equation}
\begin{gathered}
\Delta E_{\mathrm{g}}(t) = \Eg(t) - \Eg^{\mathrm{bulk}}, \\
\dEc(t_1) - \dEc(t_2) = \mathrm{EA}(t_2) - \mathrm{EA}(t_1), \qquad
\Delta E_{\mathrm{v}}(t_1) - \Delta E_{\mathrm{v}}(t_2) = \mathrm{IP}(t_1) - \mathrm{IP}(t_2) .
\end{gathered}
\label{eq:dft-partition}
\end{equation}
The bulk crystal has no vacuum, so its band edges cannot be placed on the vacuum scale directly. The partition relative
to bulk requires aligning the bulk band edges, measured from the macroscopic average potential of a bulk calculation,
with the macroscopic average inside the slab (the two-step method of \citealp{VandeWalle1987,Baldereschi1988}). This
step has not yet been done (\cref{sec:dft-pending}).

\begin{figure}[htbp]\centering
\begin{tikzpicture}[x=1cm, y=1cm, every node/.style={font=\footnotesize}, arr/.style={{Stealth[length=1.6mm]}-{Stealth[length=1.6mm]}}]
\draw[thick] (0,5.0) -- (8.6,5.0) node[right] {$E_{\mathrm{vac}}$};
% bulk
\fill[accent!25] (0.6,2.6) rectangle (2.6,3.0); \draw[thick, accent] (0.6,2.6) -- (2.6,2.6);
\fill[accent2!25] (0.6,0.6) rectangle (2.6,1.0); \draw[thick, accent2] (0.6,1.0) -- (2.6,1.0);
\node at (1.6,3.25) {CB}; \node at (1.6,0.35) {VB}; \node[font=\footnotesize\bfseries] at (1.6,5.35) {bulk};
\draw[arr] (1.0,2.6) -- node[left] {$\Eg^{\mathrm{bulk}}$} (1.0,1.0);
\draw[arr] (2.2,5.0) -- node[right] {EA} (2.2,2.6);
% slab
\fill[accent!25] (5.0,3.1) rectangle (7.0,3.5); \draw[thick, accent] (5.0,3.1) -- (7.0,3.1);
\fill[accent2!25] (5.0,0.2) rectangle (7.0,0.6); \draw[thick, accent2] (5.0,0.6) -- (7.0,0.6);
\node at (6.0,3.75) {CB}; \node at (6.0,-0.05) {VB}; \node[font=\footnotesize\bfseries] at (6.0,5.35) {thin film (slab)};
\draw[arr] (5.4,3.1) -- node[left] {$\Eg(t)$} (5.4,0.6);
\draw[arr] (6.6,5.0) -- node[right] {EA$(t)$} (6.6,3.1);
\draw[arr] (7.6,5.0) -- node[right] {IP$(t)$} (7.6,0.6);
% shifts
\draw[densely dashed, accent] (2.6,2.6) -- (4.6,2.6); \draw[densely dashed, accent2] (2.6,1.0) -- (4.6,1.0);
\draw[arr, accent] (4.4,2.6) -- node[left] {$\dEc$} (4.4,3.1); \draw[accent, thin] (4.4,3.1) -- (5.0,3.1);
\draw[arr, accent2] (4.4,1.0) -- node[left] {$\Delta E_{\mathrm{v}}$} (4.4,0.6); \draw[accent2, thin] (4.4,0.6) -- (5.0,0.6);
\end{tikzpicture}
\caption{Definitions used throughout the chapter (schematic, not to scale). Confinement in a thin film raises the
conduction-band minimum by $\dEc$ and lowers the valence-band maximum by $\Delta E_{\mathrm{v}}$, so the gap opens by
$\Delta E_{\mathrm{g}} = \dEc + \Delta E_{\mathrm{v}}$. Measured from the vacuum level, the electron affinity (EA) and
ionization potential (IP) of films of different thickness give the partition, \cref{eq:dft-partition}. The device model
so far assumes $\Delta E_{\mathrm{g}} = \dEc$.}
\label{fig:dft-levels}
\end{figure}

% =====================================================================================================================
\section{Defining the material}
\label{sec:dft-material}

\begin{dftplain}
Before anything can be computed, the computer needs a precise list of atoms: the box (cell) and the position of every
atom in it. For bulk \InO{} this comes from the measured crystal structure. A thin film is modelled as a slab: a few
layers of the crystal, cut parallel to a crystal plane, with the broken bonds at the surfaces closed by hydrogen, and
with empty space above and below. Doping is modelled by replacing one indium atom by tungsten. Every structure is
checked automatically before it is used, because a single wrong coordinate gives meaningless results.
\end{dftplain}

{\sloppy
All structures were built with the Atomic Simulation Environment \citep{Larsen2017} and checked with spglib
\citep{Togo2018} by three scripts of the package: \file{build_structure.py} for the bulk crystal, and
\file{build_structures.py} and \file{build_iwo_slab.py} for the films and the doped cells.
\Cref{tab:dft-structures} summarizes them.\par}

\subsection{Bulk \texorpdfstring{\InO}{In2O3}: the bixbyite crystal}
\label{sec:dft-bulkstruct}

Cubic \InO{} has the bixbyite structure, space group Ia$\bar{3}$ (No.~206). The conventional cubic cell, with
$a = 10.117$~\AA, contains 80 atoms on three Wyckoff sites \citep{Marezio1966}: In1 on 8b at $(\frac14,\frac14,\frac14)$,
In2 on 24d at $(u, 0, \frac14)$ with $u = -0.0336$, and O on 48e at $(0.3905, 0.1529, 0.3832)$. Each In is surrounded by
six O. The 8b site has the higher symmetry ($\bar{3}$, six equal In--O bonds); the 24d site has only a twofold axis and
three pairs of different bonds. Because the cubic cell is body-centred, the same crystal can be described by a primitive
cell of 40 atoms (16 In, 24 O), which halves the cost and is used for all bulk calculations.

The structure was generated from these crystallographic data and verified before use: space group 206 with 48
symmetry operations in the conventional cell and 24 in the primitive cell, 80 and 40 atoms, In--O bond lengths of
2.123, 2.191 and 2.208~\AA. The verification was introduced after a recorded failure: an inherited 80-atom input that
had never been checked gave only two symmetry operations, a total force of 2.86~Ry/bohr, a pressure of 765~kbar and a
negative ``gap''. A second defect followed when spglib's primitive coordinates were combined with Quantum ESPRESSO's
built-in body-centred lattice (\texttt{ibrav = 3}), whose basis vectors differ; since then all inputs give the cell
vectors explicitly (\texttt{ibrav = 0}), and pw.x reports the expected 24 operations, 18 of them with a fractional
translation (\file{PROTOCOL.md}, section 7a).

\subsection{Thin films: the slab model}
\label{sec:dft-slabstruct}

\begin{dftplain}
A plane-wave code always works with a repeating box. To describe a film, the box contains the film plus a thick layer of
empty space, so that the film does not feel its periodic copies above and below (\cref{fig:dft-slabmodel}). The
surface oxygen atoms lose bonds when the film is cut; one hydrogen atom on each of them restores a closed electronic
shell, as in a hydroxylated oxide surface.
\end{dftplain}

The slabs follow the recipe of \citet[p.~21542]{Lin2022}, which was chosen so that the results can be compared with
theirs directly:
\begin{enumerate}
\item take the conventional cell (for the 1~nm film) or a $1\times1\times2$ supercell stacked along [001] (for the 2~nm
film) of the relaxed bulk crystal;
\item remove the In layer at the cell boundary (8 cations; the layer is corrugated by about $\pm0.35$~\AA, so all
cations within 0.6~\AA{} of the boundary are removed);
\item put one H on every O that lost In neighbours, 0.97~\AA{} from the O, along the mean direction of the two removed
In--O bonds;
\item add 25~\AA{} of vacuum and centre the slab in the cell.
\end{enumerate}
The script checks the result: every one of the 24 surface O atoms has lost exactly two In neighbours; 12 H atoms sit on
each surface; the electron count closes, $3n_{\mathrm{In}} + n_{\mathrm{H}} = 2n_{\mathrm{O}}$ (In$^{3+}$, H$^+$,
O$^{2-}$), so the slab should be an insulator; no non-bonded atoms come closer than 2.61~\AA; and the space group is
Pcca (8 operations), which contains an operation that maps the top surface onto the bottom one. The slab therefore has no
dipole. Applied to the PBE-relaxed bulk ($a = 10.306$~\AA), the recipe gives In$_{24}$O$_{48}$H$_{24}$ (96 atoms,
9.875~\AA{} from the lowest to the highest H, 8.03~\AA{} between the outermost O layers) and In$_{56}$O$_{96}$H$_{24}$
(176 atoms, 20.18~\AA{} and 18.34~\AA). The films are called ``1~nm'' and ``2~nm''; Lin et al.\ quote 0.95 and 1.98~nm
for their versions \citep[p.~21540]{Lin2022}. After relaxation the 1~nm slab measures 9.30~\AA{} from H to H and
7.88~\AA{} between the outer O layers (\cref{fig:dft_structures}). For the vacuum test the 1~nm slab was also built with
15 and 35~\AA.

\begin{figure}[htbp]\centering
\begin{tikzpicture}[x=1cm, y=1cm, every node/.style={font=\footnotesize}]
\foreach \k in {0,1,2} {
  \pgfmathsetmacro{\y}{\k*2.4}
  \draw[dashed, gray] (0,\y) rectangle (3.2,\y+2.4);
  \fill[accent!35] (0,\y+0.85) rectangle (3.2,\y+1.55);
  \foreach \x in {0.2,0.6,...,3.0} { \fill[white, draw=gray] (\x,\y+1.62) circle (0.05); \fill[white, draw=gray] (\x,\y+0.78) circle (0.05); }
}
\node[right] at (3.3,3.6) {slab (film), here 1 or 2~nm};
\node[right] at (3.3,4.3) {H on the surface O atoms};
\node[right] at (3.3,2.65) {vacuum, 25~\AA};
\draw[{Stealth[length=1.5mm]}-{Stealth[length=1.5mm]}] (-0.3,2.4) -- node[left] {one cell} (-0.3,4.8);
\node[right] at (3.3,1.2) {periodic copy};
\node[right] at (3.3,6.0) {periodic copy};
\end{tikzpicture}
\caption{The slab model (schematic). The calculation repeats the cell (dashed) in all directions; along $z$ the film
alternates with 25~\AA{} of vacuum, which separates the periodic copies. Hydrogen atoms (open circles) terminate the
surface oxygen atoms.}
\label{fig:dft-slabmodel}
\end{figure}

\begin{figure}[htbp]\centering
\includegraphics[width=\linewidth]{dft_structures}
\caption{Side views ($x$--$z$ projection) of the film models. (a) Relaxed 1~nm slab; (b) 2~nm slab as built from the
relaxed bulk (relaxed geometry: \cref{sec:dft-calc-slab2}); (c) 1~nm IWO slab as built, with W (circled) on the
central cation site of the relaxed pure slab. Dashed: the periodic cell, with vacuum above and below the film. Atoms
are drawn back to front; the arrows give the distance between the outermost H atoms.}
\label{fig:dft_structures}
\end{figure}

\subsection{Tungsten doping}
\label{sec:dft-wstruct}

\begin{dftplain}
To learn where tungsten goes, one In atom of the 80-atom cell is replaced by W, once on each of the two kinds of In
site. Comparing the energies of the two relaxed cells shows which site W prefers. One W in 32 cations is 3.1\,\%,
the same order as the ``about 2\,\% W'' of the target films.
\end{dftplain}

In the conventional 80-atom cell, one In was replaced by W on an 8b site and, in a second cell, on a 24d site, giving
In$_{31}$WO$_{48}$ with W/(W+In) $= 1/32 = 3.1$\,\%. Substitution lowers the symmetry; the number of remaining operations
is 48 divided by the multiplicity of the site, which the build script checks: 6 operations (space group R$\bar{3}$) for
W on 8b, 2 (P2) for W on 24d. Relative to In$^{3+}$, W$^{6+}$ releases three electrons; in this cell of 1095~\AA$^3$
that corresponds to about $2.8\times10^{21}$\cmcu{}, far above the carrier density of the transistor channel. The cell
therefore describes the fully ionized, degenerate limit, and band-edge quantities are read from the dispersion and the
projected states, not from the position of $E_F$ alone.

For the doped film, one In of the central cation layer of the relaxed 1~nm slab was replaced by W
(\file{build_iwo_slab.py}). In the conventional-cell slabs this layer holds only 24d cations, the site preferred in
bulk (\cref{sec:dft-calc-w}). The result is In$_{23}$WO$_{48}$H$_{24}$ (W/(W+In) $= 4.2$\,\%), with W exactly at the slab
centre; the space group is P2, and its twofold operation still maps $z$ to $-z$, so the slab has no dipole. Three
electrons in a 1~nm slab of $10.34^2$~\AA$^2$ correspond to about $2.8\times10^{14}$\cmsq.

\begin{table}[htbp]\centering\footnotesize
\caption{Structures of the study. ``Electrons'' counts valence electrons of the data sets (\cref{tab:dft-pp}); ops:
symmetry operations found by spglib.}
\label{tab:dft-structures}
\begin{tabularx}{\linewidth}{@{}l>{\raggedright\arraybackslash}Xrrl>{\raggedright\arraybackslash}p{3.6cm}@{}}
\toprule
Name & Composition, cell & Atoms & Electrons & Symmetry & Purpose \\
\midrule
bulk (primitive) & In$_{16}$O$_{24}$, bcc primitive cell & 40 & 352 & Ia$\bar{3}$, 24 ops & convergence, relaxation, bands, mass \\
bulk\_conv & In$_{32}$O$_{48}$, cubic, $a = 10.306$~\AA & 80 & 704 & Ia$\bar{3}$, 48 ops & host of all other structures \\
iwo\_W8b & In$_{31}$WO$_{48}$, W on 8b & 80 & 705 & R$\bar{3}$, 6 ops & W site, valence, states \\
iwo\_W24d & In$_{31}$WO$_{48}$, W on 24d & 80 & 705 & P2, 2 ops & W site, valence, states \\
slab1\_v15/v25/v35 & In$_{24}$O$_{48}$H$_{24}$, 15/25/35~\AA{} vacuum & 96 & 624 & Pcca, 8 ops & vacuum test \\
slab1 (relaxed) & In$_{24}$O$_{48}$H$_{24}$, $10.336\times10.335\times34.87$~\AA & 96 & 624 & Pcca, 8 ops & 1~nm confinement \\
slab2\_v25 & In$_{56}$O$_{96}$H$_{24}$, $c = 45.18$~\AA & 176 & 1328 & Pcca, 8 ops & 2~nm confinement \\
slab1r\_W24d & In$_{23}$WO$_{48}$H$_{24}$ & 96 & 625 & P2, 2 ops & W in a 1~nm film \\
slab2r\_W24d & In$_{55}$WO$_{96}$H$_{24}$ & 176 & 1329 & -- & W in a 2~nm film (submitted 4 Oct.) \\
\bottomrule
\end{tabularx}
\end{table}

% =====================================================================================================================
\section{The calculations, step by step}
\label{sec:dft-calcs}

\begin{dftplain}
The work was done in stages, each building on the previous one (\cref{fig:dft-stages}). First the numerical settings
were tested on the bulk crystal; then the crystal was relaxed and its bands computed; the relaxed crystal became the
starting point for every film and doped cell. Before each stage the criteria for accepting its result were written down,
so that a result could not be accepted just because it looked good.
\end{dftplain}

\begin{figure}[htbp]\centering
\begin{tikzpicture}[node distance=4mm and 6mm, every node/.style={font=\scriptsize},
  st/.style={draw, rounded corners=2pt, align=center, fill=lightbg, text width=27mm, minimum height=9mm},
  pend/.style={st, fill=white, dashed},
  arr/.style={-{Stealth[length=1.6mm]}, thick}]
\node[st] (c) {1. Convergence\\cutoff, $k$ grid (laptop)};
\node[st, below=of c] (r) {2. Bulk vc-relax\\$a = 10.306$~\AA{} (laptop)};
\node[st, below=of r] (b) {3. Bulk bands, mass\\gap, $\mstar$ (laptop)};
\node[st, right=of c] (s) {Structures built\\from relaxed bulk};
\node[st, right=of s] (v) {4a. Vacuum test\\1~nm, 15/25/35~\AA};
\node[st, below=of v] (s1) {4b. 1~nm vc-relax\\(CPU, then GPU)};
\node[st, below=of s1] (f1) {1~nm final SCF\\EA, IP, bands};
\node[st, below=of f1] (h) {5. HSE06 check\\NC, bulk + 1~nm};
\node[st, right=of v] (w) {6. W on 8b, 24d\\relax, PDOS, spin};
\node[st, below=of w] (iw) {7. 1~nm IWO slab\\relax (GPU), final};
\node[st, below=of iw] (s2) {4c. 2~nm vc-relax\\final SCF, bands (GPU)};
\node[pend, below=of s2] (c2) {2~nm final on CPU\\(PDOS), 2~nm IWO};
\draw[arr] (c) -- (r); \draw[arr] (r) -- (b); \draw[arr] (r.east) -- ++(2mm,0) |- (s.west);
\draw[arr] (s) -- (v); \draw[arr] (v) -- (s1); \draw[arr] (s1) -- (f1); \draw[arr] (f1) -- (h);
\draw[arr] (s.north) -- ++(0,3mm) -| (w.north); \draw[arr] (w) -- (iw); \draw[arr] (s1.east) -- (iw.west);
\draw[arr] (iw) -- (s2); \draw[arr] (s2) -- (c2);
\end{tikzpicture}
\caption{Stages of the study and their dependencies. Solid: finished; dashed: running or queued at the time of
writing. Stages 1--3 ran on the author's laptop, all others on the CERN batch system.}
\label{fig:dft-stages}
\end{figure}

\subsection{Protocol and acceptance criteria}
\label{sec:dft-method}

The criteria were fixed in two protocols written before the corresponding results existed
(\file{PROTOCOL.md} of the bulk stages, 2026-09-26, and \file{PROTOCOL_CERN.md} of the CERN stages,
2026-09-27, both in \file{qe_workflow}); every later deviation is recorded there with its reason. \Cref{tab:dft-criteria} collects them.

\begin{table}[htbp]\centering\footnotesize
\caption{Acceptance criteria, fixed before the results, and their outcome.}
\label{tab:dft-criteria}
\begin{tabularx}{\linewidth}{@{}>{\raggedright\arraybackslash}p{3.0cm}>{\raggedright\arraybackslash}X>{\raggedright\arraybackslash}p{3.2cm}@{}}
\toprule
Item & Criterion & Outcome \\
\midrule
Cutoff, $k$ grid & a further increase changes the energy by $<1$~mRy/atom, the pressure by $<1$~kbar and the gap by $<10$~meV & met at 71~Ry and 2$\times$2$\times$2; 3$\times$3$\times$3 used \\
Bulk relaxation & forces $<10^{-4}$~Ry/bohr, pressure $<0.2$~kbar; lattice constant within 3\,\% of experiment & met; $+1.87$\,\% \\
$k$ grid of doped cells & 3$\times$3$\times$3 vs 4$\times$4$\times$4: energy $<1$~mRy/atom, site energy $<10$~meV & met (0.3~meV) \\
Vacuum & EA and IP change by $<10$~meV from 15 to 25 and 25 to 35~\AA & met ($\le0.1$~meV) \\
Reproduction & $\Delta E_{\mathrm{g}}$ of the 1 and 2~nm slabs within 0.1~eV of Lin et al. & met (0.04 and 0.005~eV) \\
Functional & $|\Delta E_{\mathrm{g}}^{\mathrm{HSE}}/\Delta E_{\mathrm{g}}^{\mathrm{PBE}} - 1| < 10$\,\% & met (6\,\%; 11\,\% worst case) \\
Magnetism & $|M| < 0.05\,\mu_{\mathrm{B}}$ and spin energy gain $<1$~mRy & not met; preliminary \\
Status & proxies replaced only after the $k$ and vacuum tests; ``computed'', not ``validated'' & applied \\
\bottomrule
\end{tabularx}
\end{table}

\subsection{Anatomy of an input file}
\label{sec:dft-input}

\begin{dftplain}
A calculation is described by a short text file. It says what to compute (an SCF, a relaxation, bands), how accurately
(cutoffs, thresholds), with which atoms (pseudopotentials, cell, positions) and on which $k$ points. All input files of
this study are written by scripts, never by hand, so that every job uses exactly the intended settings.
\end{dftplain}

\Cref{lst:dft-input} shows the input of the final SCF of the relaxed 1~nm slab, as written by the job
generator \file{make_jobs.py}; \cref{tab:dft-settings} lists the settings of every type of
calculation.

\begin{lstlisting}[numbers=none, caption={Input of the final SCF of the relaxed 1~nm slab (\texttt{slab1r\_final\_cpu/scf.in}); the 96 atomic positions are shortened to three.}, label={lst:dft-input}]
&CONTROL
  calculation = 'scf', prefix = 'slab', outdir = './tmp', pseudo_dir = './'
  verbosity = 'high', tprnfor = .true., tstress = .true., max_seconds = 1.0d8
/
&SYSTEM
  ibrav = 0, nat = 96, ntyp = 3, ecutwfc = 71, ecutrho = 568
  occupations = 'fixed'
  nbnd = 328
/
&ELECTRONS
  conv_thr = 1.0d-9, mixing_beta = 0.2, electron_maxstep = 150, mixing_mode = 'local-TF'
/
ATOMIC_SPECIES
  In 114.818 In.pbe-dn-kjpaw_psl.1.0.0.UPF
  O 15.999 O.pbe-n-kjpaw_psl.1.0.0.UPF
  H 1.008 H.pbe-kjpaw_psl.1.0.0.UPF
CELL_PARAMETERS angstrom
  10.3364310600 0.0000000000 0.0000000000
  0.0000000000 10.3345042100 0.0000000000
  0.0000000000 0.0000000000 34.8745389300
ATOMIC_POSITIONS crystal
  In  0.2477884238 0.2524784941 0.4261425288
  ...
  H   0.6367869492 0.3539040563 0.3667720191
K_POINTS automatic
  3 3 1 0 0 0
\end{lstlisting}

Line by line: \texttt{calculation} selects the task; \texttt{prefix} and \texttt{outdir} name the directory where the
density and orbitals are saved (later read by pp.x and by the bands run); \texttt{tprnfor} and \texttt{tstress} request
forces and stress; \texttt{max\_seconds} is overwritten by the job driver so that pw.x stops cleanly before the batch
time limit (\cref{sec:dft-driver}). \texttt{ibrav = 0} means that the cell is given explicitly in
\texttt{CELL\_PARAMETERS}; \texttt{nat} and \texttt{ntyp} count atoms and species; \texttt{ecutwfc} and \texttt{ecutrho}
are the cutoffs of \cref{sec:dft-pw}; \texttt{nbnd} is the number of bands computed, here the 312 occupied plus 16
empty ones (the empty ones give the conduction-band minimum). The \texttt{\&ELECTRONS} block sets the SCF threshold, the
mixing and the maximum number of SCF iterations. \texttt{ATOMIC\_SPECIES} gives each element its mass (IUPAC standard
atomic weight) and data set; \texttt{ATOMIC\_POSITIONS crystal} gives fractional coordinates; \texttt{K\_POINTS
automatic 3 3 1 0 0 0} requests an unshifted $3\times3\times1$ Monkhorst--Pack grid.

\begin{table}[htbp]\centering\footnotesize
\caption{Settings per type of calculation (\file{make_jobs.py}; bulk stages 1--3: \file{make_inputs.py},
\file{make_stage3.py}). All use PBE, PAW, 71/568~Ry and \texttt{conv\_thr} $= 10^{-9}$~Ry, except the hybrid check.}
\label{tab:dft-settings}
\begin{tabularx}{\linewidth}{@{}>{\raggedright\arraybackslash}p{2.5cm}>{\raggedright\arraybackslash}X>{\raggedright\arraybackslash}p{2.0cm}>{\raggedright\arraybackslash}p{4.4cm}@{}}
\toprule
Type & Occupations, mixing & $k$ points & Other settings \\
\midrule
Bulk SCF / vc-relax & fixed; $\beta = 0.4$ & 2$^3$, 3$^3$, 4$^3$ & \texttt{nbnd} 190; relax: forces $10^{-4}$~Ry/bohr, energy $10^{-5}$~Ry, pressure 0.2~kbar \\
Bulk bands & fixed; CG diagonalization & path (80), $\Gamma$ set (31) & non-SCF on the vc-relax density \\
Slab SCF & fixed; $\beta = 0.2$, local-TF; max.\ 150 iterations & 3$\times$3$\times$1 & \texttt{nbnd} = occupied + 16 \\
Slab vc-relax & as slab SCF & 3$\times$3$\times$1 & BFGS; forces $10^{-3}$~Ry/bohr, energy $10^{-4}$~Ry; \texttt{cell\_dofree = '2Dxy'}, pressure 0.5~kbar; \texttt{nbnd} = occupied + 8 \\
Slab bands & Davidson & $\Gamma$ + 10 along [100], [110] & $|k| \le 0.15$~\AA$^{-1}$ \\
Doped bulk & MV smearing 0.01~Ry; $\beta = 0.3$; max.\ 100 iterations & relax 3$^3$, SCF 4$^3$ & PDOS; bands near $\Gamma$; spin-polarized SCF 3$^3$ ($\beta = 0.2$, starting moment 0.5 on W) \\
Doped slab & MV smearing 0.01~Ry; $\beta = 0.2$, local-TF & 3$\times$3$\times$1 & relax at fixed cell; \texttt{nbnd} = occupied + 24; PDOS \\
Hybrid check & fixed; $\beta = 0.3$ (local-TF for the slab) & 3$^3$ / 3$\times$3$\times$1 & SG15 NC, 80/320~Ry; \texttt{input\_dft = 'hse'}, EXX $q$ grid 1 (bulk also 3), \texttt{ecutfock} 160~Ry \\
\bottomrule
\end{tabularx}
\end{table}

\subsection{Stage 1: convergence of the numerical settings}
\label{sec:dft-calc-conv}

\begin{dftplain}
A calculation is only meaningful if its answer does not depend on the numerical settings. The cutoff and the $k$ grid
were therefore raised step by step until energy, pressure and gap no longer changed by more than the criteria.
\end{dftplain}

\Cref{tab:dft-conv,fig:dft_convergence} summarize the tests on the 40-atom primitive cell at the experimental lattice
constant (laptop, 10 MPI processes; one SCF at 71~Ry and 3$\times$3$\times$3 took 12~min). A cutoff of 71~Ry is the lowest
that meets all three criteria against 85~Ry (0.32~mRy/atom, 0.47~kbar, 0.5~meV); 60~Ry fails the energy and pressure
criteria. The 2$\times$2$\times$2 grid already meets the criteria against 4$\times$4$\times$4; the 3$\times$3$\times$3
grid was nevertheless used for the relaxation (a recorded deviation in favour of accuracy, at almost no cost).

\begin{table}[htbp]\centering\small
\caption{Convergence of bulk \InO{} at $a = 10.117$~\AA{} (40-atom primitive cell, PAW, PBE). Energies relative to the
best-converged run of each series.}
\label{tab:dft-conv}
\begin{tabular}{lrrrr}
\toprule
Run & $E$ (Ry) & $\Delta E$/atom (mRy) & $P$ (kbar) & KS gap (eV) \\
\midrule
50~Ry, 3$\times$3$\times$3 & $-7573.24296524$ & 3.17 & 84.28 & 1.1897 \\
60~Ry, 3$\times$3$\times$3 & $-7573.32959316$ & 1.00 & 89.87 & 1.1894 \\
71~Ry, 3$\times$3$\times$3 & $-7573.35691861$ & 0.32 & 96.57 & 1.1897 \\
85~Ry, 3$\times$3$\times$3 & $-7573.36970781$ & 0 & 97.04 & 1.1902 \\
71~Ry, 2$\times$2$\times$2 & $-7573.35524852$ & $+0.044$ vs 4$\times$4$\times$4 & 96.27 & 1.1886 \\
71~Ry, 4$\times$4$\times$4 & $-7573.35701116$ & reference & 96.61 & 1.1898 \\
\bottomrule
\end{tabular}
\end{table}

\begin{figure}[htbp]\centering
\includegraphics[width=\linewidth]{dft_convergence}
\caption{Convergence of total energy, pressure and Kohn--Sham gap of bulk \InO{} with the plane-wave cutoff (left,
3$\times$3$\times$3) and the $k$-point grid (right, 71~Ry). The grey band marks the energy criterion of 1~mRy per atom;
the dotted lines on the right mark the gap criterion of 10~meV.}
\label{fig:dft_convergence}
\end{figure}

\subsection{Stage 2: the relaxed bulk crystal}
\label{sec:dft-calc-relax}
\label{sec:dft-bulk}

At the experimental lattice constant PBE gives a pressure of +97~kbar: the crystal ``wants'' to expand, the expected
behaviour of this functional. The variable-cell relaxation (71~Ry, 3$\times$3$\times$3, BFGS, 1~h~47~min on the laptop)
converged in 11 BFGS steps to a primitive volume of 547.333~\AA$^3$, i.e.\ $a = 10.306$~\AA, 1.87\,\% above the
experimental 10.117~\AA{} \citep{Marezio1966} and within 0.06\,\% of the PBE value of 10.30~\AA{} reported by
\citet[p.~21542]{Lin2022}. The final SCF with the plane-wave basis of the new cell gives $P = 0.41$~kbar, inside the
1~kbar criterion. Space group Ia$\bar{3}$ is preserved; the In--O bonds become 2.215~\AA{} (8b) and 2.167/2.240/2.256~\AA{}
(24d). All later structures are built from this relaxed cell.

\subsection{Stage 3: bulk band structure and effective mass}
\label{sec:dft-calc-bands}

Two non-SCF runs use the density of the relaxed crystal: one along the high-symmetry path
$\Gamma$--H--N--$\Gamma$--P--H$|$P--N of the body-centred cubic Brillouin zone \citep{Setyawan2010} (80 $k$ points,
about 15 per $2\pi/a$; 11~h~16~min), and one on $\Gamma$ plus 10 points along each of [100], [110] and [111] up to
$|k| = 0.15$~\AA$^{-1}$ (31 points; 4~h~21~min). The first attempt with Davidson diagonalization aborted after 9 of 16
$k$ points with unconverged eigenvalues; both runs therefore use conjugate-gradient diagonalization.

\Cref{fig:dft_bulk_bands} shows the bands. The conduction-band minimum lies at $\Gamma$; the valence band is very flat,
with its maximum on $\Gamma$--H at $k = (0,0,0.196)\,2\pi/a$, only 4.4~meV above the highest valence state at $\Gamma$.
The fundamental gap is 0.889~eV and the direct gap at $\Gamma$ 0.893~eV, against 0.94~eV in
\citet[p.~21542]{Lin2022}; the bulk gap of 3.05~eV adopted in the device model is, as expected, strongly underestimated
by PBE (\cref{sec:dft-xc}). The bulk reference for all confinement shifts is the fundamental gap, 0.8887~eV.

\begin{figure}[htbp]\centering
\includegraphics[width=0.75\linewidth]{dft_bulk_bands}
\caption{PBE band structure of bulk \InO{} at the relaxed lattice constant, energies relative to the valence-band
maximum. Occupied bands blue, empty bands red. Path after \citet{Setyawan2010}.}
\label{fig:dft_bulk_bands}
\end{figure}

\Cref{fig:dft_mass_fit}(a) shows how the mass is obtained. Within $|k| \le 0.05$~\AA$^{-1}$ the parabolic fit
(\cref{eq:dft-mass}) gives $\mstar = 0.1588$, 0.1587 and 0.1587~$m_0$ along [100], [110] and [111]: an isotropic
$\mstar = 0.159\,m_0$, as expected for the s-like conduction-band minimum of a cubic crystal. Further from $\Gamma$ the
points fall below the parabola: the band is non-parabolic. The Kane fit (\cref{eq:dft-kane}) over $|k| \le 0.153$~\AA$^{-1}$
gives $\alpha = 0.63$, 0.57 and 0.55~eV$^{-1}$. For comparison, \citet[p.~21540]{Lin2022} report $0.17\,m_0$ with PBE, and
the device model uses the measured $0.208\,m_0$ \citep[p.~225102-12]{Stokey2021}. PBE underestimates the mass because
the gap is too small (the conduction-band curvature is coupled to the gap); as for the gap, only differences between
films and bulk computed in the same setup are transferred to the device model.

\begin{figure}[htbp]\centering
\includegraphics[width=\linewidth]{dft_mass_fit}
\caption{How the effective mass is obtained. Symbols: conduction-band energies from the non-SCF runs, relative to the
minimum at $\Gamma$; dashed: parabola from the fit over $|k| \le 0.05$~\AA$^{-1}$ (shaded). (a) Bulk, three directions;
(b) relaxed 1~nm slab, two in-plane directions. The departure from the parabola at larger $k$ is the non-parabolicity.}
\label{fig:dft_mass_fit}
\end{figure}

\subsection{Stage 4a: vacuum thickness}
\label{sec:dft-calc-vac}

\begin{dftplain}
The vacuum layer must be thick enough that the film does not ``see'' its periodic copies. This was tested by computing
the same film with 15, 25 and 35~\AA{} of vacuum: if EA and IP do not change, the vacuum is sufficient.
\end{dftplain}

The unrelaxed 1~nm slab was computed with the three vacuum thicknesses (16 CPUs each; the 35~\AA{} SCF took 1~h~48~min
and 77 iterations). EA and IP change by at most 0.1~meV and the vacuum plateau of the potential is flat to 0.01~meV
(\cref{tab:dft-vac}), far inside the 10~meV criterion; the symmetric, dipole-free slab converges fast. The 25~\AA{} of
\citet[p.~21542]{Lin2022} is therefore sufficient and is used for all slabs.

\begin{table}[htbp]\centering\small
\caption{Vacuum convergence of the unrelaxed 1~nm slab (In$_{24}$O$_{48}$H$_{24}$).}
\label{tab:dft-vac}
\begin{tabular}{lrrrr}
\toprule
Vacuum & Gap (eV) & EA (eV) & IP (eV) & SCF iterations \\
\midrule
15~\AA & 0.9586 & 3.4331 & 4.3917 & 46 \\
25~\AA & 0.9586 & 3.4332 & 4.3918 & 64 \\
35~\AA & 0.9586 & 3.4332 & 4.3918 & 77 \\
\bottomrule
\end{tabular}
\end{table}

\subsection{Stage 4b: relaxation and confinement of the 1~nm film}
\label{sec:dft-calc-slab1}
\label{sec:dft-slab1}

\begin{dftplain}
A freshly cut film is not in equilibrium: the surface atoms and the added hydrogens are pulled into new positions. The
relaxation lets them move. For the 1~nm film this changed the result completely: the gap of the unrelaxed film was only
0.07~eV above the bulk gap, the relaxed one 0.90~eV.
\end{dftplain}

The relaxation (all atoms and the in-plane cell) started on 16 CPU cores; after 13 BFGS steps and the successful GPU
benchmark (\cref{sec:dft-gpu}) it was continued on one A100 GPU from the latest geometry, where it converged after 54
further steps (15~h~31~min). \Cref{fig:dft_relaxation_history} shows the 72 SCF cycles: the total energy fell by
7.05~eV (73~meV per atom), the residual force by a factor of 40, and the gap rose from 0.959 to 1.789~eV, most of it in
the first ten steps. The final in-plane cell is
$10.336\times10.335$~\AA{} (bulk 10.306; Lin et al.\ 10.26 and 10.36~\AA, \citealp[p.~21542]{Lin2022}). The total
residual force of 0.0061~Ry/bohr is the norm over all 96 atoms; BFGS ended with the criteria as Quantum ESPRESSO applies
them.

\begin{figure}[htbp]\centering
\includegraphics[width=\linewidth]{dft_relaxation_history}
\caption{Relaxation of the 1~nm slab: total energy relative to the final value, Kohn--Sham gap and total force at every
SCF cycle (one per geometry). Dotted: continuation of the CPU run on a GPU from its latest geometry. The gap of the
relaxed film is the one used.}
\label{fig:dft_relaxation_history}
\end{figure}

The confinement shift is
\begin{equation}
\Delta E_{\mathrm{g}}(\SI{0.99}{nm}) = 1.7885 - 0.8887 = \SI{0.900}{eV},
\label{eq:dft-deg1}
\end{equation}
against $+0.94$~eV for the 0.95~nm slab of \citet[p.~21542]{Lin2022}. The in-plane mass is $0.281\,m_0$ along [100] and
$0.282\,m_0$ along [110] (\cref{fig:dft_mass_fit}b), an increment of $+0.122\,m_0$ over the bulk value, against
$+0.13\,m_0$ in \citet[p.~21540]{Lin2022}. Both agree within the pre-registered tolerance: the published confinement of a
1~nm film is reproduced independently, with a different pseudopotential type and cutoff (PAW at 71~Ry versus
norm-conserving at 140~Ry, \citealp[p.~21542]{Lin2022}).

\begin{dftplain}
\Cref{fig:dft_lin_parity} puts every quantity that both studies computed side by side. A point on the vertical line
would mean perfect agreement; the grey band is $\pm5$\,\%. All our energies and masses lie a few per cent below theirs,
and our lattice constants agree to 0.2\,\%.
\end{dftplain}

The lattice constants agree within 0.2\,\%. Gaps and masses are uniformly 4--7\,\% lower than those of
\citet[pp.~21540, 21542]{Lin2022}: bulk gap $-5.5$\,\%, 1~nm gap $-4.9$\,\%, bulk mass $-6.6$\,\%, 1~nm mass $-6.3$\,\%.
The quantities transferred to the device model, the differences between film and bulk, show the same small offset
($\Delta E_{\mathrm{g}}$ $-4.3$\,\%, $\Delta\mstar$ $-5.8$\,\%), well inside the pre-registered tolerance of 0.1~eV for the
shift. An offset of the same sign and size for every quantity points to a systematic difference between the two
numerical setups (pseudopotential type and cutoff) rather than to a difference between the films.

\begin{figure}[htbp]\centering
\includegraphics[width=0.92\linewidth]{dft_lin_parity}
\caption{This work against the PBE calculations of \citet[pp.~21540, 21542]{Lin2022}, which used the same slab recipe:
relative difference of each quantity (dots) with both values written out. Grey: $\pm5$\,\%; green: the pre-registered
tolerance of the confinement shift ($\pm0.1$~eV; at 2~nm this is $\pm30$\,\% and the bar is clipped). For the in-plane
lattice constant of the 1~nm film, Lin et al.\ give 10.26 and 10.36~\AA; their mean is compared with ours. Below the
dotted line: the 2~nm film (\cref{sec:dft-calc-slab2}); no criterion was set for the masses. The increments of Lin et
al.\ are differences of their rounded values (1.27 $-$ 0.94~eV and 0.23 $-$ 0.17\,$m_0$).}
\label{fig:dft_lin_parity}
\end{figure}

\begin{dftplain}
\Cref{fig:dft_levels} shows the allowed electron energies at the centre of the Brillouin zone ($\Gamma$) as horizontal
lines, for the bulk crystal and for the 1~nm film before and after relaxation. Below zero are the filled (valence)
levels, above the gap the empty (conduction) ones. Two things change in the film: the gap opens, and new, separate
conduction levels appear above the lowest one. They are the signature of confinement: an electron that can move freely
through the crystal can only take certain energies when it is trapped between two surfaces, like a wave on a string of
fixed length.
\end{dftplain}

At $\Gamma$ the gap is 0.893~eV in bulk, 1.015~eV in the film as built (whose valence-band maximum lies away from
$\Gamma$, so its fundamental gap is 0.959~eV) and 1.788~eV after relaxation. In the bulk crystal the next conduction
level at $\Gamma$ lies 4.28~eV above the minimum; in the film there are further levels 1.57 and 2.43~eV (as built) and
1.50 and 2.45~eV (relaxed) above the lowest one. This is consistent with the conduction band splitting into separate
sub-bands when the motion across the film is confined. The valence levels are denser in the film because its cell
contains more atoms.

\begin{figure}[htbp]\centering
\includegraphics[width=0.72\linewidth]{dft_levels}
\caption{Kohn--Sham levels at $\Gamma$ near the gap, relative to the highest occupied level at $\Gamma$: bulk \InO{}
(40-atom cell), the 1~nm film as built, and the relaxed 1~nm and 2~nm films (PBE). Blue: occupied; red: empty. Arrows:
gap at $\Gamma$.}
\label{fig:dft_levels}
\end{figure}

The pp.x step of the GPU job failed to read the density written by the GPU build (\cref{sec:dft-gpu}), so the final
geometry was recomputed on 16 CPU cores (SCF 2~h~19~min, 47 iterations; bands 4~h~38~min), which reproduced the gap and
mass to the last printed digit and gave the potential of \cref{fig:dft_planar_potential}: EA $=3.93$~eV and IP
$=5.72$~eV (PBE). \Cref{fig:dft_confinement_vs_thickness} places the result next to the literature points and the
thickness laws of the device model.

\begin{figure}[htbp]\centering
\includegraphics[width=0.8\linewidth]{dft_planar_potential}
\caption{Planar (grey) and macroscopic (blue) averages of the electrostatic potential of the relaxed 1~nm slab along the
film normal, relative to the vacuum level, with the conduction- and valence-band edges (PBE). The flat regions near
$z = 0$ and $z = c$ are the vacuum.}
\label{fig:dft_planar_potential}
\end{figure}

\begin{dftplain}
\Cref{fig:dft_alignment} draws the band edges of three 1~nm films on one common scale: the energy of an electron at
rest outside the film (the vacuum level) is zero, and every edge is placed below it. This is the computed version of
the schematic \cref{fig:dft-levels}.
\end{dftplain}

On the vacuum scale, relaxation lowers the conduction-band minimum of the 1~nm film from $-3.43$ to $-3.93$~eV and the
valence-band maximum from $-4.39$ to $-5.72$~eV. Both edges move down, by different amounts. Besides the film interior,
relaxation changes the surface O--H bonds, and any change of their electric dipole shifts the potential step between
film and vacuum; absolute EA and IP of a single film therefore contain a contribution of the surface termination, which
this calculation does not separate. Only differences between films with the
same termination (here the relaxed 1 and 2~nm films) isolate the confinement shift of each band edge
(\cref{eq:dft-partition}). The third column is discussed in \cref{sec:dft-calc-iwoslab}.

\begin{figure}[htbp]\centering
\includegraphics[width=0.78\linewidth]{dft_alignment}
\caption{Band edges of the films on the vacuum scale, computed from the SCF eigenvalues and the planar-averaged
potential (PBE): 1~nm film as built and relaxed, relaxed 2~nm film, and 1~nm film with one W on the central 24d site.
For the pure films, the band edges are taken
over the whole $k$ grid, as in \cref{tab:dft-vac}. For the W-doped film, the levels at $\Gamma$ are shown: the top
valence level, the bottom of the host conduction band (red; dark: filled up to $E_F$), the W-5d-derived level (purple)
and $E_F$ (dashed). That column is preliminary.}
\label{fig:dft_alignment}
\end{figure}

\begin{figure}[htbp]\centering
\includegraphics[width=\linewidth]{dft_confinement_vs_thickness}
\caption{Confinement shift (left) and mass increment (right) versus film thickness: thickness laws of the device model,
PBE slab values of \citet[pp.~21540, 21542]{Lin2022}, the conduction-band shift of \citet[p.~504]{Si2021} (corundum
\InO{} on \AlO, PBE), and this work (PBE, and the HSE06-corrected range of \cref{sec:dft-hse}, drawn slightly offset in
thickness for visibility). Thickness of this work: H--H distance of the slab as built (0.99 and 2.02~nm). The
HSE06 correction measured at 1~nm is applied unchanged at 2~nm (an assumption, not a calculation).}
\label{fig:dft_confinement_vs_thickness}
\end{figure}

\subsection{Stage 5: is PBE adequate for the shift? The hybrid-functional check}
\label{sec:dft-hse}

\begin{dftplain}
PBE puts the gap far too low. If it also got the \emph{change} of the gap wrong, the confinement values would be
useless. The check repeats the bulk and the 1~nm film with the much more accurate HSE06 functional and compares the two
shifts. They differ by only 6\,\%.
\end{dftplain}

The check uses HSE06 single points at the PBE geometries. It was first run with PAW; after more than 4.5~h on an H100
GPU at 0\,\% utilization (and more than 5.5~h on 16 CPU cores before) the first exact-exchange step had not finished,
the PAW exact exchange of the GPU build evidently running on the host. Following a recorded deviation the check was
moved to SG15 norm-conserving pseudopotentials \citep{Hamann2013,Schlipf2015} (In, O, H at 80~Ry, a cutoff verified
against 100~Ry: gap change 0.1~meV), with PBE and HSE06 computed in the same setup and the ratio of the shifts
transferred. The exact-exchange cutoff \texttt{ecutfock} was set to twice the wave-function cutoff after the default ran
the GPU out of memory for the slab (the output reports 40.6~GB for the exchange arrays at this setting); for the bulk
this changes the HSE06 gap by less than 0.1~meV.

A hybrid calculation is a loop within a loop: pw.x first converges a PBE SCF (34 iterations for the slab), then
repeatedly fixes the exchange operator, converges the density in it (23, 14, 9, 5, 3, 3, 1 iterations), and updates the
operator, until the estimated exchange error falls from 0.013~Ry to $1.4\times10^{-9}$~Ry. The slab took 4.0~h and the
bulk 12~min on one H100 NVL GPU.

\begin{table}[htbp]\centering\small
\caption{PBE versus HSE06 (norm-conserving, same geometries): Kohn--Sham gaps at $\Gamma$ and the confinement shift.}
\label{tab:dft-hse}
\begin{tabular}{lrrr}
\toprule
 & Bulk gap (eV) & 1~nm slab gap (eV) & $\Delta E_{\mathrm{g}}$ (eV) \\
\midrule
PBE & 0.9204 & 1.8068 & 0.8864 \\
HSE06 ($q = 1$) & 2.1080 & 3.0475 & 0.9395 \\
\midrule
Ratio HSE06/PBE & & & 1.060 \\
\bottomrule
\end{tabular}
\end{table}

HSE06 raises both gaps by about 1.2~eV but changes the shift by only 6\,\% (\cref{tab:dft-hse,fig:dft_pbe_vs_hse}): the
ratio 1.060 is inside the pre-registered 10\,\% criterion, so the PBE confinement is accepted. The reduced
exact-exchange $q$ grid is the dominant numerical uncertainty: with $q = 3\times3\times3$ the bulk HSE06 gap is 42~meV
lower (2.0662~eV); if the slab gap did not shift by the same amount, the ratio would be 1.107, so it lies between 1.06
and 1.11. The norm-conserving PBE shift (0.886~eV) agrees with the PAW value (0.900~eV) within 1.6\,\%, which justifies
applying the ratio to the PAW result: the HSE-corrected 1~nm shift is 0.95--1.00~eV. The semi-local functional,
inadequate for the absolute gap, is adequate for the confinement \emph{shift} to within about 10\,\%, which supports the
slab-minus-bulk construction of the device model (\cref{sec:der-confinement}) with evidence rather than assumption.

\begin{figure}[htbp]\centering
\includegraphics[width=0.62\linewidth]{dft_pbe_vs_hse}
\caption{Kohn--Sham gaps at $\Gamma$ of bulk \InO{} and of the relaxed 1~nm slab with PBE and HSE06 (norm-conserving,
same geometries) and the resulting confinement shifts.}
\label{fig:dft_pbe_vs_hse}
\end{figure}

\subsection{Stage 6: tungsten in bulk \texorpdfstring{\InO}{In2O3}}
\label{sec:dft-calc-w}
\label{sec:dft-w}

\begin{dftplain}
Three questions were asked about W: which of the two In sites it takes, which charge state (valence) it has, and
whether its extra electrons become free carriers. The site follows from comparing energies, the valence from the
length of the W--O bonds, and the fate of the electrons from where their states sit in energy and on which atoms.
\end{dftplain}

Each doped cell was relaxed at the PBE lattice constant (3$\times$3$\times$3, 32 CPU cores; 10 and 24~h), followed by an
SCF on 4$\times$4$\times$4 (30 irreducible points for W on 24d), a PDOS, bands near $\Gamma$ and a spin-polarized SCF
(\cref{tab:dft-settings}). \Cref{tab:dft-w} gives the results.

\begin{table}[htbp]\centering\small
\caption{W in the 80-atom \InO{} cell (In$_{31}$WO$_{48}$, 3.1\,\% W).}
\label{tab:dft-w}
\begin{tabular}{@{}>{\raggedright\arraybackslash}p{5.4cm}>{\raggedright\arraybackslash}p{3.0cm}>{\raggedright\arraybackslash}p{3.6cm}@{}}
\toprule
Quantity & W on 8b & W on 24d \\
\midrule
$E$, relaxed, 3$\times$3$\times$3 (Ry) & $-15496.3476466$ & $-15496.3666037$ \\
$E$, SCF 4$\times$4$\times$4 (Ry) & $-15496.34766084$ & $-15496.36660308$ \\
W--O bonds (\AA) & 6 $\times$ 1.972 & 1.892, 1.902, 2.152 (each $\times$2; mean 1.982) \\
$E_F$ (eV), 4$\times$4$\times$4 & 10.758 & 10.794 \\
Spin-polarized: moment ($\mu_{\mathrm{B}}$/cell); $E_{\mathrm{spin}} - E$ & 0.74; $-1.67$~mRy & 0.62; $-1.43$~mRy \\
W 5d weight of the three bands at and above $E_F$ (first $k$ point) & 0.50, 0.50, 0.58 & 0.55, 0.48, 0.30 \\
\bottomrule
\end{tabular}
\end{table}

\paragraph{Site.} The energy difference
\begin{equation}
E(\mathrm{8b}) - E(\mathrm{24d}) = +0.258~\mathrm{eV}\ (3\times3\times3),\qquad +0.2577~\mathrm{eV}\ (4\times4\times4),
\label{eq:dft-site}
\end{equation}
is converged to 0.3~meV against the 10~meV criterion: W prefers the 24d site. With the multiplicity ratio of three, the
equilibrium occupation ratio is $N_{24d}/N_{8b} = 3\exp(\Delta E/k_{\mathrm{B}}T)$, about 430:1 at 600~K.

\paragraph{Valence.} The W--O bonds are about 0.2~\AA{} shorter than In--O. Bond lengths expected from the sums of the
Shannon radii \citep{Shannon1976} (O$^{2-}$ 1.38~\AA, six-coordinated cations) are 1.98~\AA{} for W$^{6+}$, 2.00 for
W$^{5+}$, 2.04 for W$^{4+}$ and 2.18 for In$^{3+}$; the computed means of 1.97--1.98~\AA{} identify W$^{6+}$. Relative to
In$^{3+}$, each W can therefore release three electrons.

\paragraph{The electrons of W (preliminary).} The projected densities of states (\cref{fig:dft_w_pdos}) show that the
states at and just above $E_F$ carry about half of their weight on W 5d orbitals (\cref{tab:dft-w}), while the
valence-band top carries 1.5\,\%. A spin-polarized calculation finds weak moments of 0.6--0.7~$\mu_{\mathrm{B}}$ per cell,
1.4--1.7~mRy below the non-magnetic state; this does not meet the pre-registered non-magnetic criterion and is treated
as preliminary, because coarse sampling of a degenerate electron gas can produce spurious moments. At the first $k$
point the host conduction band (band 352, 67\,\% In 5s) lies 1.76~eV below $E_F$, and $E_F$ lies in the
W-5d-derived band (band 353, 54\,\% W 5d). Within PBE and for one W per cell (3.1\,\% of the cations,
$5.5\times10^{21}$ donated electrons per cm$^3$), the donated electrons thus fill the host conduction band up to a
W-5d-derived level that is resonant in the conduction band, about 1.8~eV above its bottom, and partly occupy that level.
How the three electrons divide between the two is not determined from the available data.

\subsection{Stage 7: tungsten in the 1~nm film}
\label{sec:dft-calc-iwoslab}

The W-doped 1~nm slab (\cref{sec:dft-wstruct}) was relaxed at the cell of the pure relaxed slab on one H100 NVL GPU
(27 BFGS steps, 4~h~29~min; $E = -12238.28910896$~Ry); the final SCF, potential and PDOS ran on 16 CPU cores
(2~h~55~min), and its bands run aborted, so no mass is available. The W--O bonds are 1.880, 1.907 and 2.212~\AA{}
(each twice).

\paragraph{Which band is which.} The doped slab has one occupied band \emph{fewer} below the gap than the pure slab:
the In data set treats the filled 4d shell as valence (five bands per In), the W data set the 5s and 5p shells (four
bands), so replacing one In by W removes one band. The projections at $\Gamma$ confirm this: band 311 is the top of
the O 2p valence band ($-2.366$~eV), band 312 ($-0.647$~eV) is the bottom of the host conduction band (In 5s 50\,\%,
O 2s 21\,\%), band 313 (0.576~eV) carries 56\,\% W 5d weight and band 314 (0.766~eV) is the next host conduction
level. A first version of this analysis took band 312 as the valence-band top, as in the pure slab; it therefore
reported a W level 0.19~eV \emph{below} the conduction band and an unexplained upward shift of 1.1--1.65~eV of all
levels on the vacuum scale. Both were artefacts of the band count and were corrected on 5 October 2026.

\paragraph{Result (PBE, one W position, preliminary).} $E_F$ lies 1.23~eV above the bottom of the host conduction band
and 5~meV above band 313. The host band is filled: band 312 lies below $E_F$ at every $k$ point and holds two of the
three donated electrons; about one electron occupies the W-5d-derived level, which is resonant in the conduction band,
1.22~eV above its bottom. The gap between the valence-band top and the host conduction-band bottom is 1.719~eV, 0.069~eV
less than in the pure film. On the vacuum scale (\cref{fig:dft_alignment}, fourth column) the doped film has
EA = 4.07~eV and IP = 5.79~eV, 0.14 and 0.07~eV larger than the pure film: the levels of the two films agree within
0.14~eV.

\begin{figure}[htbp]\centering
\includegraphics[width=\linewidth]{dft_w_pdos}
\caption{Total density of states (grey, scaled) and W 5d projected density of states (red, scaled) relative to $E_F$:
bulk 80-atom cell with W on 24d (left) and 1~nm slab with W on the central 24d site (right). PBE, preliminary.}
\label{fig:dft_w_pdos}
\end{figure}

\paragraph{Consequence for the device model (inference).} Full activation of about 2\,\% W would give of the order of
$10^{21}$\cmcu{} electrons, orders of magnitude more than the channels of transistors that switch off can carry. In
both cells the W-5d-derived level is resonant in the conduction band, 1.2~eV (1~nm film) to 1.8~eV (bulk) above its
bottom, and takes up electrons only because these cells fill the band up to it (about $3\times10^{21}$ and
$5.5\times10^{21}$ donated electrons per cm$^3$; in the 1~nm film, about one of the three electrons). The first result
for the 2~nm film with 1.8\,\% W (last SCF of its relaxation; projections pending) places $E_F$ 1.24~eV above the
conduction-band bottom and the flat bands that are presumably W-derived 0.07~eV above $E_F$, i.e.\ just empty. At the
lower electron densities of working channels, PBE thus predicts that W releases its electrons into the host band;
trapping in W 5d states cannot by itself explain why the channels switch off, and compensation by other defects or
charge at the interfaces must account for the difference. A hybrid functional can lower localized d levels, so this
conclusion is preliminary. The calculations do not determine the effective donor density and support the treatment of
the channel donor density as a fitted, film-specific parameter. The W content of the target films is known only as
``about 2\,\% W'' without a stated basis (\cref{ch:device}); the W fractions of the cells (3.1 and 4.2\,\% of the
cations) are of the same order.

\subsection{Stage 4c: the 2~nm film}
\label{sec:dft-calc-slab2}

\begin{dftplain}
The 2~nm film is the thickness of the thinnest measured device, so it is the most important single calculation for the
device model. It answers three questions: how much does the gap open at 2~nm, how much heavier do the electrons
become, and is the change in the conduction band or in the valence band? The answers: the gap opens by 0.34~eV (Lin et
al.\ found 0.33~eV), the mass rises by about 30\,\%, and, going from 2 to 1~nm, the whole additional opening is in the
conduction band.
\end{dftplain}

The 176-atom slab is the largest calculation of the study: 1328 electrons, a $45.18$~\AA{} cell, and about 24~GB of
memory per $k$-point pool on CPUs (it was held at a 48~GB limit with four pools). Its variable-cell relaxation ran on one
H100 NVL GPU with the settings of the 1~nm slab. A first run converged after 68 BFGS steps (2 October 2026) but was
lost: its final SCF, with more bands than the relaxation, needed 66~GB of memory against the 32~GB requested, the job
was held by the batch system, and no checkpoint of it had reached EOS (\cref{tab:dft-failures}). The relaxation was
resubmitted from the same starting geometry with 100~GB and performed the final SCF, the planar average (with the
density converted for pp.x, \cref{sec:dft-gpu}) and the bands on the GPU itself.

\paragraph{Relaxation.} The second run converged after 70 BFGS steps (72 SCF cycles, 41.5~h; final enthalpy
$-27035.66548$~Ry, total residual force 0.0072~Ry/bohr). The relaxed in-plane cell is $10.318\times10.309$~\AA{} (1~nm
film: $10.336\times10.335$; bulk 10.306~\AA). The film contracts in thickness: the H--H distance falls from 20.18 to
19.62~\AA{} and the O--O distance from 18.34 to 18.21~\AA{} (1~nm film: H--H 9.87 to 9.30~\AA). The final SCF took
56~min and the bands run 8~h~22~min on the GPU; the whole job ran for 50.8~h.

\paragraph{Gap and mass.} The gap of the relaxed film is 1.2239~eV, at $\Gamma$ like that of the 1~nm film, so
\begin{equation}
\Delta E_{\mathrm{g}}(\SI{2.02}{nm}) = 1.2239 - 0.8887 = \SI{0.335}{eV},
\label{eq:dft-deg2}
\end{equation}
against $+0.33$~eV for the 1.98~nm slab of \citet[p.~21542]{Lin2022} (difference 0.005~eV, inside the pre-registered
0.1~eV). The in-plane mass is $0.206\,m_0$ along both [100] and [110] (0.2064 and 0.2059), an increment of
$+0.047\,m_0$ over bulk. Lin et al.\ give $0.23\,m_0$ \citep[p.~21540]{Lin2022}: our mass is 10\,\% lower, against
6\,\% lower at 1~nm, so the increment ($+0.047$ versus $+0.06\,m_0$) is 20\,\% lower; for the increment the rounding of
their values (two digits) is of the same order as the difference. At $\Gamma$, the next empty level lies 0.72~eV above
the conduction-band minimum (1.50~eV in the 1~nm film, \cref{fig:dft_levels}): the confinement levels move closer
together as the film gets thicker.

\paragraph{Comparison with the device model.} The thickness laws of the device model (\cref{eq:dec-law}), fitted to
the published proxies, give 0.348~eV and $0.049\,m_0$ at 2.02~nm; the computed values are 4\,\% and 3\,\% lower
(\cref{fig:dft_confinement_vs_thickness}). At 0.99~nm the laws give 0.933~eV and $0.133\,m_0$, against the computed 0.900~eV
and $0.122\,m_0$. A refit to the computed points will therefore change the laws by a few per cent.

\paragraph{Where the change sits.} On the vacuum scale (\cref{fig:dft_alignment}) the relaxed 2~nm film has EA = 4.495~eV
and IP = 5.719~eV, with a flat vacuum plateau. Both films have the same termination and the same read-out, so their
difference partitions the change of the gap between 1 and 2~nm:
\begin{equation}
\begin{gathered}
\dEc(1) - \dEc(2) = \mathrm{EA}(2) - \mathrm{EA}(1) = +0.567~\mathrm{eV}, \\
\Delta E_{\mathrm{v}}(1) - \Delta E_{\mathrm{v}}(2) = \mathrm{IP}(1) - \mathrm{IP}(2) = -0.002~\mathrm{eV},
\end{gathered}
\end{equation}
whose sum, $+0.565$~eV, equals $\Delta E_{\mathrm{g}}(1) - \Delta E_{\mathrm{g}}(2)$. Between 1 and 2~nm the whole change of
the gap is thus in the conduction band, and the ionization potential is independent of thickness within 2~meV. Two
checks are open before this is used. First, a top valence level that does not move at all may belong to the surface
(O--H) rather than to the film interior; a projection of that state onto the atoms will decide this. Second, the shares
relative to bulk, $\dEc(t)$ and $\Delta E_{\mathrm{v}}(t)$ themselves, need the bulk potential at the slab's in-plane
lattice constant (two-step alignment, \cref{sec:dft-readout}). The CPU cross-check of the final SCF (which also gives
the projected density of states) and the 2~nm IWO slab were submitted on 4 October 2026.

% =====================================================================================================================
\section{How the calculations were run: the workflow}
\label{sec:dft-compute}

\begin{dftplain}
The bulk stages fit on a laptop; the films do not. A single 2~nm calculation needs tens of gigabytes of memory and days
of computing. The film and doping calculations therefore ran on CERN's batch system, which distributes jobs over
thousands of processor cores and a pool of GPUs. A collection of scripts prepares the jobs, sends them, watches them,
repairs common failures, collects the results and analyses them, so that the calculations could run day and night
without manual intervention. Logging in remained a manual step done by the author: no script stores or types a
password.
\end{dftplain}

\subsection{Computers and software}
\label{sec:dft-cern}

Stages 1--3 ran on the author's laptop under Windows Subsystem for Linux with Quantum ESPRESSO 7.5 (conda-forge build)
and 10 MPI processes. All later calculations ran on the CERN batch system, which uses HTCondor \citep{Thain2005}: jobs
are submitted from the login service lxplus, wait in a queue and are matched to free worker nodes; each job declares its
CPUs, memory, disk, GPUs and a ``flavour'' that sets its maximum wall time (here 8~h, 1, 3 or 7 days). Files are
exchanged through EOS, CERN's disk storage. Access used a Kerberos ticket and an SSH connection that the author opened
and authenticated manually; the scripts reuse this connection and never handle credentials.

The CPU jobs run Quantum ESPRESSO 7.5 from a relocatable copy of the conda environment stored on EOS and unpacked on the
worker. The GPU jobs run NVIDIA's GPU build of Quantum ESPRESSO 7.3.1 \citep{NvidiaQE} from a container image on EOS,
executed with Apptainer \citep{Kurtzer2017}. Structure building and analysis use Python with ASE, spglib, NumPy
\citep{Harris2020} and Matplotlib \citep{Hunter2007}.

\subsection{Anatomy of a job}
\label{sec:dft-driver}

\file{make_jobs.py} turns the list of calculations into a job directory (\file{jobs/gen_<name>/}) containing: one
sub-directory per job with its input files (\file{*.in}, templates \file{*.tmpl} that receive the relaxed geometry
later, and \file{steps.sh}, the list of steps); the pseudopotentials; \file{jobs.txt} (one line per job: name, CPUs,
memory, disk, flavour, time limit, GPUs); \file{jobs_info.json} (atoms, electrons, irreducible $k$ points); the submit
description \file{jobs.sub}; the worker script \file{driver.sh}; and two helpers, \file{qeio.py} (geometry transfer,
convergence test, Fermi level) and \file{dat2h5.py} (density conversion, \cref{sec:dft-gpu}).

On the worker, \file{driver.sh}:
\begin{enumerate}
\item unpacks Quantum ESPRESSO from EOS (and, for GPU jobs, copies the container image);
\item forces one thread per MPI process (the variables \texttt{OMP\_NUM\_THREADS},
\texttt{OPENBLAS\_NUM\_THREADS} and \texttt{MKL\_NUM\_THREADS} are set to 1) and restricts MPI to shared-memory communication on one node;
\item runs the steps of \file{steps.sh}. Each pw.x step receives \texttt{max\_seconds} = time left $-$ 30~min, so that
it stops cleanly and writes its results before the batch limit; a guard stops a run whose output exceeds 2~GB (a sign
of a looping error);
\item after every step copies the outputs so far to EOS (a checkpoint), and at the end packs all outputs into one
archive that HTCondor returns.
\end{enumerate}
A typical slab job (\file{steps.sh} of \file{slab1r_final_cpu}) is three lines: \texttt{run\_pw scf 4}, \texttt{run\_ppavg
slab \$JOB 4.8833}, \texttt{run\_pw bands 4}. The number after \texttt{run\_pw} is the number of $k$-point pools: the
MPI processes are split into groups that each handle some of the $k$ points, which is efficient but multiplies the
memory. \file{make_jobs.py} chooses the largest number of pools that divides the CPUs, does not exceed the irreducible $k$
points and leaves at least four processes per pool; the 2~nm slab is limited to two pools by memory. A relaxation step
is followed by \texttt{qeio.py converged} (did it finish?) and \texttt{qeio.py newgeom}, which copies the final cell and
positions into the templates of the following SCF and bands runs.

\subsection{GPU acceleration}
\label{sec:dft-gpu}

Before any GPU result was used, a finished CPU calculation (the unrelaxed 1~nm slab) was repeated on one H100 NVL GPU
with identical input: total energies $-11888.29271202$ and $-11888.29271155$~Ry ($5\times10^{-7}$~Ry apart), the same
gap and the same 46 SCF iterations, in 9~min~16~s against 1~h~16~min on 16 CPU cores, a factor of 8.2
(\cref{fig:dft_computing}). The difference of QE versions (7.3.1 on GPU, 7.5 on CPU) has no effect at the protocol
tolerances. GPU jobs request one GPU and 8 CPU cores; one MPI process drives the GPU and the remaining cores serve as
OpenMP threads; only H100 NVL and H200 nodes are accepted (A100 nodes ran about two times slower).

Two incompatibilities had to be solved. First, the PAW exact exchange of the GPU build runs on the host
(\cref{sec:dft-hse}). Second, the GPU build saves the density as a plain Fortran file (\file{charge-density.dat}), whereas
the CPU build used for pp.x and projwfc.x, compiled with HDF5, reads only \file{charge-density.hdf5}; pp.x aborted on the
first GPU save, and the band edges of GPU runs were lost. The converter \file{dat2h5.py} rewrites the density in the HDF5
layout of QE 7.5 (header attributes, Miller indices with the reciprocal vectors, the complex density as interleaved real
numbers), using the h5py library of CERN's software distribution on the worker. It was tested by a round trip: a QE 7.5
density converted to the plain format and back gave identical HDF5 headers and data and a byte-identical planar average
from pp.x and average.x. The 2~nm final calculation is the first production use.

\subsection{The automated workflow}
\label{sec:dft-flow}

All operations are commands of one documented script, \file{qe_workflow/dft_flow.sh} (\cref{tab:dft-flow}), also available
as tasks in Visual Studio Code. \Cref{fig:dft-workflow} shows the flow of data.

\begin{table}[htbp]\centering\footnotesize
\caption{Commands of \file{dft_flow.sh}.}
\label{tab:dft-flow}
\begin{tabularx}{\linewidth}{@{}>{\ttfamily}p{3.4cm}X@{}}
\toprule
\normalfont Command & What it does \\
\midrule
login, check, keepalive & print the login command for the author to type; test the SSH session; keep it alive and renew the Kerberos ticket every 6~h, warning 24~h before the ticket can no longer be renewed \\
build & build and verify all structures from the relaxed bulk \\
generate, plan, submit & write job directories; show the resources each job would get; upload and submit, skipping jobs already queued or finished \\
status, resources, hw & queue and live progress (step, SCF iterations, BFGS steps, energy); free CPU slots and GPUs by model; hardware under each job \\
fetch, analyze & download finished results (and checkpoints of jobs that died); collect all numbers into \file{summary.json} \\
auto & every 5--10~min: status, resources, fetch, watchdog, chains, analysis \\
dashboard & live web page of the whole study \\
local & run one job on the laptop with the same driver \\
\bottomrule
\end{tabularx}
\end{table}

\begin{figure}[htbp]\centering
\begin{tikzpicture}[node distance=5mm and 7mm, every node/.style={font=\scriptsize},
  box/.style={draw, rounded corners=2pt, align=center, fill=lightbg, text width=27mm, minimum height=9mm},
  arr/.style={-{Stealth[length=1.6mm]}, thick}]
\node[box] (b) at (0,0) {1. Laptop (WSL)\\build structures,\\generate jobs};
\node[box] (l) at (4.4,0) {2. lxplus: submit\\(SSH session opened\\by the author)};
\node[box] (q) at (8.8,0) {3. HTCondor queue\\matches CPU / GPU};
\node[box] (w) at (8.8,-2.0) {4. Worker node\\\file{driver.sh} runs\\the QE steps};
\node[box] (e) at (4.4,-2.0) {5. EOS storage\\QE software, GPU image,\\checkpoints, results};
\node[box] (f) at (0,-2.0) {6. Laptop: fetch,\\analyse, \file{summary.json}};
\node[box] (r) at (0,-3.9) {7. Figures, report,\\dashboard};
\node[box, text width=36mm] (a) at (8.8,-3.9) {\texttt{auto} loop: watchdog,\\race/upgrade, continuation,\\chains};
\draw[arr] (b) -- (l); \draw[arr] (l) -- (q); \draw[arr] (q) -- (w);
\draw[arr] ([yshift=2mm]e.east) -- node[above, font=\tiny] {software} ([yshift=2mm]w.west);
\draw[arr] ([yshift=-2mm]w.west) -- node[below, font=\tiny] {outputs} ([yshift=-2mm]e.east);
\draw[arr] (e) -- node[above, font=\tiny] {fetch} (f);
\draw[arr] (f) -- (r);
\draw[arr, dashed] (a) -- node[right, font=\tiny] {supervises} (w);\end{tikzpicture}
\caption{Data flow of the workflow. Jobs are generated on the laptop, submitted through the author's lxplus session,
matched to CPU or GPU workers by HTCondor, and run by \file{driver.sh}, which takes the software from EOS and writes
checkpoints there. Results return to the laptop for analysis. The \texttt{auto} loop supervises the queue.}
\label{fig:dft-workflow}
\end{figure}

\paragraph{Allocation.} \texttt{submit} decides each job's resources from its size and the free resources at that
moment: films with 90 or more atoms get one H100 NVL or H200 GPU with 8 CPU cores if one is free, otherwise 16 CPU
cores; smaller jobs get 16 CPU cores (32 only when more than 20 such slots are free, because jobs of 32 or more cores go
to a small pool of 10 slots and waited 6~h on the first day).

\paragraph{Supervision.} The \texttt{auto} loop applies fixed rules:
\begin{itemize}
\item \emph{race and upgrade}: copies of the same calculation (a GPU job and its CPU twin) form a group; the first copy
that runs removes queued copies of equal or lower rank, and a GPU copy that starts removes a running CPU copy;
\item \emph{idle}: a GPU job queued for more than 3~h gets a CPU twin; a CPU job queued for more than 3~h is resubmitted
with a smaller request;
\item \emph{stuck}: a job that runs for 60~min in an SCF step without a single SCF iteration is resubmitted (at most
twice); a job without progress for 30~min is reported;
\item \emph{memory}: a job held for exceeding its memory is released with 1.5 times the memory;
\item \emph{continuation}: a relaxation that ends unfinished (time limit, eviction) is resubmitted from its last
geometry as \texttt{JOB\_c}$N$ (at most five times);
\item \emph{chains}: a command attached to a relaxation runs once its finished result has been fetched; the 2~nm
relaxation is chained to its final runs and to the 2~nm IWO slab.
\end{itemize}
These rules were tested with deliberate fault-injection jobs (``canaries'' that simulate an unfinished relaxation, a
stuck start and a race), which exposed and fixed several rule errors before the rules acted on production jobs.

\subsection{Failures and what they taught}
\label{sec:dft-failures}

\Cref{tab:dft-failures} lists the failures in the order they occurred. Each is recorded in the protocols and logs; none
affected a reported number, because every result comes from a run that finished cleanly and passed its checks.

\begin{table}[htbp]\centering\footnotesize
\caption{Failures, causes and fixes.}
\label{tab:dft-failures}
\begin{tabularx}{\linewidth}{@{}>{\raggedright\arraybackslash}p{3.8cm}>{\raggedright\arraybackslash}X>{\raggedright\arraybackslash}p{4.0cm}@{}}
\toprule
Symptom & Cause & Fix \\
\midrule
2 symmetry operations, 765~kbar, negative gap & inherited, never-checked bulk input & rebuilt from crystallographic data, automatic checks \\
again 2 operations & spglib coordinates with QE's \texttt{ibrav = 3} basis & cell given explicitly (\texttt{ibrav = 0}) \\
32-core jobs stopped after 18~s & MPI launcher counted physical cores, slots expose hardware threads & \texttt{--map-by :OVERSUBSCRIBE} \\
2~nm slab held at 95.8~GB & four $k$-point pools, about 24~GB each & at most two pools \\
bulk bands aborted (Davidson) & unconverged eigenvalues in a non-SCF run & CG for bulk bands; Davidson kept for slabs (CG took 13.7~h instead of 1.5~h) \\
jobs 17--30 times too slow & batch system set thread counts to the core count; 65 threads per process & one thread per process, all jobs resubmitted \\
large jobs waiting for hours & 32+ cores go to a 10-slot pool & 16-core CPU jobs, 8-core GPU jobs \\
2~nm unrelaxed SCF lost at 24~h & killed at the limit in the bands step & checkpoint to EOS after every step \\
GPU job wrote 324~GB of errors & MPI tried InfiniBand on GPU nodes & transport restricted to shared memory and CUDA; 2~GB output guard \\
hybrid functional at 0\,\% GPU & PAW exact exchange on the host in the GPU build & norm-conserving pseudopotentials \\
GPU out of memory (hybrid, slab) & default exact-exchange cutoff & \texttt{ecutfock} = 2 $\times$ \texttt{ecutwfc} \\
pp.x could not read GPU saves & plain vs HDF5 density format & \file{dat2h5.py} converter \\
converged 2~nm relaxation lost & final SCF needed 66~GB (32~GB requested); job held, no checkpoint on EOS (cause not determined) & memory from measured demand; no automatic release after $>1$~h of running; re-run \\
workflow rule errors & logic errors found by canary jobs & fixed before production use \\
\bottomrule
\end{tabularx}
\end{table}

\subsection{Computing resources}
\label{sec:dft-resources}

The usage to date is about 8\,300 weighted CPU core-hours and 85--90 GPU-hours at CERN, of which about 3\,000
core-hours and 50 GPU-hours were lost to the failures above, and about 190 core-hours on the laptop
(\cref{fig:dft_computing}). Run on the laptop, the useful work would take about six weeks of continuous computing, and
the 2~nm slab and the hybrid slab calculation (30--100~GB of memory) could not run at all. \Cref{tab:dft-timing}
gives the wall times of representative steps.

\begin{table}[htbp]\centering\footnotesize
\caption{Wall times of representative steps (from the output files and driver logs).}
\label{tab:dft-timing}
\begin{tabular}{@{}lllr@{}}
\toprule
Calculation & Size & Hardware & Wall time \\
\midrule
bulk SCF, 71~Ry, 3$\times$3$\times$3 & 40 atoms & laptop, 10 MPI & 12~min \\
bulk vc-relax (11 BFGS steps) & 40 atoms & laptop, 10 MPI & 1~h~47~min \\
bulk bands, path (80 $k$, CG) & 40 atoms & laptop, 10 MPI & 11~h~16~min \\
1~nm slab SCF, unrelaxed & 96 atoms & 16 CPU cores & 1~h~48~min \\
same, GPU benchmark & 96 atoms & 1 H100 NVL & 9~min \\
1~nm vc-relax, 54 BFGS steps & 96 atoms & 1 A100 & 15~h~31~min \\
1~nm final SCF / bands & 96 atoms & 16 CPU cores & 2~h~19~min / 4~h~38~min \\
W on 24d: relax / SCF 4$^3$ / bands & 80 atoms & 32 CPU cores & 23~h~48~min / 3~h~3~min / 18~h~36~min \\
HSE06 bulk / 1~nm slab & 40 / 96 atoms & 1 H100 NVL & 12~min / 4.0~h \\
IWO 1~nm relax (27 BFGS steps) & 96 atoms & 1 H100 NVL & 4~h~29~min \\
\bottomrule
\end{tabular}
\end{table}

\begin{figure}[htbp]\centering
\includegraphics[width=\linewidth]{dft_computing}
\caption{Left: CERN CPU usage (approximate weighted core-hours), productive and lost to faults that were fixed. Right:
GPU benchmark, the same 96-atom SCF on 16 CPU cores and on one H100 NVL.}
\label{fig:dft_computing}
\end{figure}

\begin{dftplain}
\Cref{fig:dft_timeline} is the diary of the CERN work: one row per calculation, coloured by what the calculation was
doing at each moment (waiting in the queue, running on CPUs, running on a GPU). It shows the long CPU phase of the
first two days, the move to GPUs once they were validated, and the many short checks that ran in parallel on
30 September.
\end{dftplain}

The timeline combines three records: the exact start and end times written by \file{driver.sh} for every job that
returned output; the queue snapshots of the monitor script every 30~min from 27 September 09:02Z; and the 5--10~min
status records of the workflow from 28 September 14:20Z. Runs of the first submission that were removed without output
before 09:02Z on 27 September (the thread problem, \cref{tab:dft-failures}) are not shown, apart from the failed
18-second starts. White gaps within a row mean that no copy of that calculation was recorded as queued or running (between a
removal and a resubmission, or while the monitoring was interrupted).

\begin{figure}[htbp]\centering
\includegraphics[width=\linewidth]{dft_timeline}
\caption{Timeline of the CERN calculations, 27 September to 1 October 2026 (UTC). Blue: running on CPUs; red: running
on a GPU; grey: queued; amber: held; hatched: failed at the first step. Dotted lines: the thread fix, the GPU
validation, and the move of the hybrid check to norm-conserving pseudopotentials. Several copies of a calculation (CPU
twin, GPU variants, continuations) share one row.}
\label{fig:dft_timeline}
\end{figure}

% =====================================================================================================================
\section{From output files to numbers and figures}
\label{sec:dft-analysis}

\begin{dftplain}
A finished calculation leaves text files of thousands of lines. Small Python scripts read the few lines that matter
(energies, eigenvalues, forces, the averaged potential), compute the derived quantities with the formulas of
\cref{sec:dft-readout}, and draw the figures. Nothing is copied by hand, so every number in this chapter can be
regenerated from the raw outputs by running the scripts again.
\end{dftplain}

\subsection{Reading the outputs}
\label{sec:dft-parsing}

\Cref{tab:dft-parsing} lists each quantity, the line of the output it comes from and the script that reads it. The
eigenvalues of non-SCF runs are read block by block (one block per $k$ point, \cref{lst:dft-eig}); $k$ is converted from
units of $2\pi/a_{\mathrm{lat}}$ to \AA$^{-1}$ with the lattice parameter printed in the same output. This matters after
a vc-relax: pw.x keeps the initial \texttt{alat} as its unit, so the largest $k$ of the bulk $\Gamma$ set is 0.153 and
not 0.15~\AA$^{-1}$.

\begin{table}[htbp]\centering\footnotesize
\caption{Where each quantity comes from.}
\label{tab:dft-parsing}
\begin{tabularx}{\linewidth}{@{}>{\raggedright\arraybackslash}p{3.0cm}>{\raggedright\arraybackslash}X>{\raggedright\arraybackslash}p{3.4cm}@{}}
\toprule
Quantity & Source in the output & Script \\
\midrule
total energy & line starting with ``\texttt{!}'' (last one); ``Final energy/enthalpy'' after a relaxation & \file{collect_results.py}, \file{make_dft_figures.py} \\
pressure, force & ``\texttt{P=}'' in the stress block; ``\texttt{Total force =}'' & same \\
VBM, CBM, gap (SCF) & ``highest occupied, lowest unoccupied level (ev)'' & \file{analyze_slabs.py} \\
bulk fundamental gap & minimum of band 177 minus maximum of band 176 over the path run & \file{analyze_path.py}, \file{collect_results.py} \\
$\mstar$, $\alpha$ & eigenvalue blocks of the $\Gamma$-set bands run, fits \cref{eq:dft-mass,eq:dft-kane} & \file{analyze_mass.py}, \file{collect_results.py} \\
$E_{\mathrm{vac}}$, EA, IP & \file{<job>_avg.dat} of average.x: mean of the planar average within 2~\AA{} of the cell boundary; its spread as a flatness check & \file{analyze_slabs.py} \\
$E_F$ & ``the Fermi energy is'' & \file{qeio.py fermi} \\
W--O bonds & final geometry, neighbours within 2.6~\AA{} (ASE) & \file{analyze_iwo.py} \\
W 5d share & projwfc.x output and \file{pdos/*(W)_wfc#*(d)} files & \file{analyze_iwo_slab.py} \\
site energy & $[E(\mathrm{8b}) - E(\mathrm{24d})] \times 13.6057$~eV/Ry & \file{collect_results.py} \\
HSE ratio & gaps of the four NC runs & \file{collect_results.py} \\
\bottomrule
\end{tabularx}
\end{table}

\begin{lstlisting}[language=Python, caption={Reading eigenvalues and fitting the effective mass (\texttt{make\_dft\_figures.py}, shortened).}, label={lst:dft-eig}]
def eig(path, nocc, extra=8):
    t = Path(path).read_text(errors='replace')
    alat = float(re.search(r'lattice parameter \(alat\)\s+=\s+([\d.]+)', t).group(1)) * BOHR
    bl = re.findall(r'k =\s*([-\d. ]+?)\s*\(\s*\d+ PWs\)\s+bands \(ev\):\s+(.*?)(?=...)', t, re.S)
    ks = np.array([...]) * 2 * np.pi / alat          # k in 1/A
    ev = np.array([...])                              # eigenvalues (eV), bands 1..nocc+extra
    return ks, ev

k = np.linalg.norm(ks[idx], axis=1); e = ev[idx, nocc] - ev[0, nocc]   # CB relative to its minimum
sel = k <= 0.0501                                     # fit window |k| <= 0.05 1/A
c2 = np.polyfit(k[sel] ** 2, e[sel], 1)[0]            # slope of E against k^2
m_star = H2M / c2                                     # H2M = hbar^2/(2 m0) = 3.80998 eV A^2
\end{lstlisting}

All numbers are collected by \file{collect_results.py} into the file \file{summary.json} of the results
directory, which feeds the live dashboard; values that are missing stay empty and are never estimated. Every value used in this chapter is also entered,
with its date and context, in \file{RESULTS_LOG.md}.

\subsection{Drawing the figures}
\label{sec:dft-figures}

All figures of this chapter are produced by one script of the report, \file{tools/make_dft_figures.py}
(Matplotlib, Times font, PDF output), run with \file{python tools/make_dft_figures.py}. Each figure is one function
that reads the output files directly (\cref{tab:dft-figures}). The plotting choices follow the purpose of each figure:
energies are always shown relative to a physical reference (the VBM, the CBM, $E_F$ or $E_{\mathrm{vac}}$), never as raw
eigenvalues, because absolute Kohn--Sham energies have no meaning; acceptance criteria are drawn as shaded bands so that
a reader sees at once whether a point passes; literature values are open symbols and values of this work filled.

\begin{table}[htbp]\centering\footnotesize
\caption{The figure functions of \file{make_dft_figures.py}: data and processing.}
\label{tab:dft-figures}
\begin{tabularx}{\linewidth}{@{}>{\raggedright\arraybackslash}p{3.2cm}>{\raggedright\arraybackslash}p{4.1cm}>{\raggedright\arraybackslash}X@{}}
\toprule
Figure & Data read & Processing and plotting \\
\midrule
\cref{fig:dft_structures} & structure files (\file{*.json}) of the relaxed and built slabs & $x$--$z$ projection; atoms drawn back to front by $y$, radius per element; cell outline; H--H distance \\
\cref{fig:dft_convergence} & \file{scf_ecut*_k*.out} & differences to the most converged run, per atom (40 atoms) in mRy; pressure and gap differences; criteria shaded \\
\cref{fig:dft_bulk_bands} & \file{bands_path.out}, \file{bands_path.in} & distance along the path from successive $|\Delta k|$, no distance across the H$|$P jump; tick positions from the points per segment of the input; energies minus the VBM \\
\cref{fig:dft_mass_fit} & \file{bands_gamma.out}, slab \file{bands.out} & $|k|$ from $\Gamma$, $E - E_{\mathrm{CBM}}$; least-squares slope against $k^2$ in the window; parabola drawn over the full range \\
\cref{fig:dft_relaxation_history} & vc-relax outputs of the CPU run and the GPU continuation & energy (log scale), gap and total force for each SCF cycle; the switch marked \\
\cref{fig:dft_planar_potential} & \file{slab1r_final_cpu_avg.dat}, \file{scf.out} & columns $z$ (bohr), planar and macroscopic average (Ry) converted to \AA{} and eV; $E_{\mathrm{vac}}$ from the vacuum window; band edges as lines \\
\cref{fig:dft_confinement_vs_thickness} & gap of \file{slab1r_final_cpu}; literature values; law coefficients & $\Delta E_{\mathrm{g}} = \Eg(t) - 0.8887$~eV; HSE-corrected point $\times1.085\pm0.025$ (midpoint and half-range of 1.06--1.11); laws of the device model as curves \\
\cref{fig:dft_pbe_vs_hse} & four NC \file{scf.out} & gaps as bars, shifts and ratio in the panel \\
\cref{fig:dft_w_pdos} & \file{pdos_tot}, W 5d PDOS files, Fermi energy & energies relative to $E_F$; each curve scaled to its maximum (shapes, not absolute weights) \\
\cref{fig:dft_computing} & resource accounting of the logs; benchmark wall times & stacked bar; bar chart of the benchmark \\
\cref{fig:dft_lin_parity} & bulk and slab outputs (as above); values of Lin et al.\ with page & relative difference per quantity; $\pm5$\,\% band; tolerance of the shift \\
\cref{fig:dft_levels} & eigenvalue block at $\Gamma$ of \file{bands_gamma.out} and of the two slab \file{scf.out} & levels as horizontal lines relative to the highest occupied level at $\Gamma$; occupied/empty by band index \\
\cref{fig:dft_alignment} & \file{scf.out} and \file{*_avg.dat} of three slabs & $E_{\mathrm{vac}}$ from the vacuum window; band edges (pure films: whole $k$ grid; doped film: $\Gamma$ levels 312--314 and $E_F$) placed below it \\
\cref{fig:dft_timeline} & \file{monitor.log}, \file{status_history.jsonl}, all \file{*_driver.log} & samples converted to bars of one sampling interval per row (copies merged); driver-log intervals drawn exactly \\
\bottomrule
\end{tabularx}
\end{table}

% =====================================================================================================================
\section{How the gap opening is shared: alignment to bulk}
\label{sec:dft-partition}

\begin{dftplain}
When the film becomes thin, the gap opens. For the transistor only the upward move of the conduction band matters,
because that is where the electrons travel. To split the opening into its conduction-band and valence-band parts, the
energies of the film and of the bulk crystal must be placed on one common scale. We used the average electrostatic
potential deep inside the film, which there looks like bulk, as the common reference.
\end{dftplain}

\paragraph{Method.} The alignment is done in two steps
(\file{qe_workflow/bulk_In2O3_protocol/run_bulk_potential.sh}, \file{bulk_alignment.py}).

\emph{Step 1, bulk.} The cell average of the local electrostatic potential (\texttt{pp.x}, \texttt{plot\_num} 11) of
the relaxed bulk crystal is $\langle V\rangle = 4.1687$~eV. Relative to it, the band edges of the bulk band structure
are $E_v - \langle V\rangle = 3.9098$~eV and $E_c - \langle V\rangle = 4.7985$~eV.

\emph{Step 2, slab.} The macroscopic average of the same potential is taken at the centre of each slab.

\paragraph{Averaging window.} The planar-average period of bixbyite along [001] is $a/2$, not the $a/4$ used for the
vacuum levels.
\begin{itemize}
\item With an $a/4$ window the slab interior oscillates by 177--205~meV. The result disagrees with the
electron-affinity route by about 80~meV, so it was rejected.
\item The $a/2$ window gives an interior plateau flat within 8~meV at 2~nm (1~meV with the double $a/2\ast a/4$
average).
\end{itemize}

\paragraph{The 1~nm film.} It has almost no bulk-like interior, so its values follow from the 2~nm result plus the
same-termination electron-affinity differences ($+0.567$ and $-0.002$~eV). Between 1 and 2~nm the two routes agree
within 15--30~meV.

\begin{table}[htbp]\centering\small
\caption{Band-edge shifts relative to bulk (PBE). $\dEc - \Delta E_{\mathrm{v}} = \Delta E_{\mathrm{g}}$ in every case.}
\label{tab:dft-partition}
\begin{tabular}{llccc}\toprule
Film & Route & $\dEc$ (eV) & $\Delta E_{\mathrm{v}}$ (eV) & conduction-band share\\\midrule
2~nm & two-step, $a/2$ window & $+0.281$ & $-0.054$ & 0.84\\
2~nm & double average & $+0.287$ & $-0.049$ & 0.86\\
1~nm & two-step, $a/2$ window & $+0.819$ & $-0.081$ & 0.91\\
1~nm & 2~nm value + electron-affinity difference (adopted) & $+0.848$ & $-0.056$ & 0.94\\\bottomrule
\end{tabular}
\end{table}

\paragraph{Caveats.}
\begin{itemize}
\item The in-plane strain of the slabs relative to bulk ($+0.29$\,\% at 1~nm, $+0.07$\,\% at 2~nm) is not corrected;
its effect is likely a few meV at 2~nm.
\item The values are PBE.
\end{itemize}

\paragraph{The band edges are not surface states.} An orbital projection of the relaxed 2~nm film (\texttt{projwfc.x},
\file{qe_workflow/cern_htcondor/surface_state_check.py}) checked this.
\begin{itemize}
\item The highest occupied state carries 0.2\,\% of its weight on the surface OH groups, which make up 27\,\% of the
atoms. Its weight peaks at the centre of the film, like the lowest confined state of a box.
\item The lowest empty state carries 1.4\,\% on the OH groups.
\item A deeper valence state used as a reference carries 14\,\%.
\end{itemize}
Both band edges are therefore confined, bulk-like states, and the split of \cref{tab:dft-partition} describes them.

\paragraph{Consequence for the device model.} Model V1 assumed the whole opening in the conduction band. The computed
share is 0.84--0.94. At 2~nm the conduction band therefore moves up by 0.28~eV, not 0.35~eV. The device model takes
these values over in \cref{ch:v2}.

% =====================================================================================================================
\section{Pending results and changes to the device model}
\label{sec:dft-pending}

\begin{table}[htbp]\centering\small
\caption{Status of the first-principles inputs of the device model.}
\label{tab:dft-status}
\begin{tabularx}{\linewidth}{@{}>{\raggedright\arraybackslash}p{4.2cm}>{\raggedright\arraybackslash}X>{\raggedright\arraybackslash}p{3.3cm}@{}}
\toprule
Input & This work & Status \\
\midrule
$\Delta E_{\mathrm{g}}$ at about 1~nm & +0.90~eV (PBE); 0.95--1.00~eV (HSE-corrected) & computed; reproduces Lin et al. \\
$\Delta\mstar$ at about 1~nm & $+0.122\,m_0$ & computed \\
$\Delta E_{\mathrm{g}}$ at about 2~nm & +0.335~eV (PBE); 0.36--0.37~eV (HSE ratio of 1~nm assumed) & computed; reproduces Lin et al. \\
$\Delta\mstar$ at about 2~nm & $+0.047\,m_0$ & computed \\
$\dEc/\Delta E_{\mathrm{v}}$ partition & relative to bulk: 2~nm $\dEc=+0.281$, $\Delta E_{\mathrm{v}}=-0.054$~eV (share 0.84); 1~nm $+0.848$, $-0.056$~eV (share 0.94); between 1 and 2~nm all in $\dEc$ (\cref{sec:dft-partition}) & computed; band edges bulk-like (not surface states) at 2~nm \\
Films of 0.75 and 0.51~nm & Lin recipe with two and one cation planes (In$_{16}$O$_{36}$H$_{24}$, In$_8$O$_{24}$H$_{24}$; H to H 0.75 and 0.51~nm before relaxation; \file{structures_v2/build_thin_slabs.py}, verified to reproduce the 1~nm slab exactly). The one-plane film has the composition of indium hydroxide, In(OH)$_3$, and marks the limit of the recipe & relaxing on CERN (5 Oct.\ 2026) \\
W site, valence & 24d, W$^{6+}$ & computed, converged \\
W 5d level, moments & resonant in the conduction band, 1.2--1.8~eV above its bottom; occupied only at very high filling (\cref{sec:dft-calc-w,sec:dft-calc-iwoslab}) & preliminary (PBE); corrected 5 Oct.\ 2026 \\
W effect on confinement & \textit{[to be completed: 2~nm IWO slab; 1~nm IWO slab mass]} & pending \\
\bottomrule
\end{tabularx}
\end{table}

The 1 and 2~nm points are now computed. They lie 3--8\,\% below the V1 laws (\cref{sec:dft-calc-slab2}). With the
partition relative to bulk (\cref{sec:dft-partition}), the conduction-band shift at 2~nm is 22\,\% below the V1 law.

\Cref{ch:v2} refits the confinement and mass laws to the computed points (model V2) and gives them the provenance label
of this work. It also calibrates the device model again and tests it on the held-out 6.3~nm film.

% =====================================================================================================================
\section{Limitations}
\label{sec:dft-limits}

\begin{caveat}
The films are idealized crystalline (001) slabs with hydrogen-passivated surfaces in vacuum; the device channels are
sputtered, largely amorphous or nanocrystalline IWO on \AlO{} with an air-exposed top surface. The results describe
trends and provide physically grounded inputs; they do not predict the measured devices directly. PBE is adequate for
the confinement shift within about 10\,\% (\cref{sec:dft-hse}) but not for absolute gaps, masses or the position of W 5d
states. The W cells contain one W at 3.1--4.2\,\%; other positions, concentrations, complexes with oxygen defects and
charge states were not computed. The hybrid check uses a reduced exact-exchange $q$ grid, whose effect is bounded but
not eliminated. The relaxations stop at forces of $10^{-3}$~Ry/bohr (a recorded deviation from the $10^{-4}$ of Lin et
al.).
\end{caveat}

\begin{keyresult}
An independent PAW-PBE calculation reproduces the published confinement of 1 and 2~nm \InO{} films ($\Delta
E_{\mathrm{g}} = +0.90$ and $+0.335$~eV, $\Delta\mstar = +0.122$ and $+0.047\,m_0$, against $+0.94$ and $+0.33$~eV and
$+0.13$ and $+0.06\,m_0$). Between 1 and 2~nm the change of the gap lies entirely in the conduction band, and the band
edges are bulk-like, not surface states. A hybrid-functional check shows that the
semi-local functional gives the confinement shift to within 6--11\,\%, so the slab-minus-bulk construction of the device
model rests on evidence. Tungsten substitutes on the 24d site as W$^{6+}$; within PBE its donated electrons
fill the host conduction band, and a W-5d-derived level resonant 1.2--1.8~eV above the band bottom takes up electrons
only at very high filling (preliminary). Every structure, setting, run and number is traceable to scripts and output files,
and the workflow that produced them is documented. The computed points lie 3--8\,\% below the V1 thickness laws, and 84--94\,\%
of the gap opening is in the conduction band (V1 assumed 100\,\%); the device model takes these values over in
\cref{ch:v2}. The W effect on confinement is pending.
\end{keyresult}
